\chapter{初等数论函数与素数计数函数}\label{ch:elementary_number_theoretic_functions}

% ============================================================
\section{引言}\label{sec:intro_elementary_functions}

研究素数分布（以及后文将系统展开的黎曼 $\zeta$ 函数与黎曼猜想），
需要一组定义在实数或正整数上的初等中间函数。
它们分别承担如下工作：用 $\pi$、$J$、$\theta$、$\psi$ 记录素数及其对数加权的个数；
用 $\mu$、$\Lambda$ 作算术反演与卷积；用 Mertens 函数 $M(x)=\sum_{n\le x}\mu(n)$ 记录 $\mu$ 的部分和；
用偏移对数积分 $\Li$ 给出素数计数的光滑比较主项
（完整 $\Li$/$\li$/$\Ei$ 专章见第~\ref{ch:logint_Ei}~章）。
在讨论 $\zeta$ 的解析延拓、零点与显式公式之前，须先清楚定义算术侧函数并建立彼此关系。
它们与 $\zeta$ 的 Dirichlet 级数联系（$\sum\mu(n)n^{-s}=1/\zeta(s)$、
$\sum\Lambda(n)n^{-s}=-\zeta'/\zeta$）在第~\ref{ch:riemann_zeta_intro}~章于 $\operatorname{Re}s>1$ 上建立。

各节给出具体函数的定义、基本恒等式、函数图形和与其他函数的互换关系。
章末用进度图标出截至本章已建立的主要对象与关系。

本章对象按定义方式可分为三类（后文 Perron 与显式公式主要针对第一类的和函数）：

\begin{center}
\renewcommand{\arraystretch}{1.35}
\begin{tabular}{|c|c|p{6.2cm}|}
\hline
\textbf{类型} & \textbf{成员} & \textbf{定义与计算特征} \\
\hline
算术阶梯 / 部分和
  & $\pi,\,J,\,\theta,\,\psi,\,M$
  & 由素数计数或系数累加规定；一般无初等闭式通项，
    算具体值须枚举（或等价筛法） \\
\hline
算术系数（定义在正整数 $n$ 上）
  & $\mu,\,\Lambda$
  & 按 $n$ 的素因子分解分情况定义；给定分解后可算，
    其部分和即上列 $M,\,\psi$ 等 \\
\hline
解析比较主项
  & $\Li$
  & 实积分 $\int_2^x\mathrm{d}t/\log t$；用作 $\pi$ 的光滑近似。
    $\li$、$\Ei$ 与复 $\li(x^{\rho})$ 见第~\ref{ch:logint_Ei}~章 \\
\hline
\end{tabular}
\end{center}

\noindent
因此 $\Li$ 与 $\pi$ 等阶梯不是同一类对象：前者是解析比较函数，后者是本章要立名的算术阶梯。

% ============================================================
\section{素数计数函数 \texorpdfstring{$\pi(x)$}{pi(x)}}\label{sec:prime_counting}

本节给出素数个数的标准计数函数：定义、小范围数值与阶梯图。
后文的 $J$、$\theta$、$\psi$ 及素数定理均以 $\pi$ 为参照对象之一。

\subsection{定义}
 
\begin{definition}{}{}
对实数 $x \geq 0$，\textbf{素数计数函数} $\pi(x)$ 定义为不超过 $x$ 的素数个数：
\[
  \pi(x) \;=\; \sum_{\substack{p \leq x \\ p\,\text{素数}}} 1.
\]
\end{definition}
 

黎曼论文《论小于给定数值的素数个数》中「小于给定数值的素数个数」，正是 $\pi(x)$ 的定义。
上式中的求和目前只是计数记号，而不是可直接计算的封闭形式：
在给定 $x$ 时，仍须逐个检验不超过 $x$ 的整数是否为素数并计数。
本书围绕 $\pi(x)$ 的解析表示与渐近公式展开推导与证明。

$\pi(x)$ 的图像是一条阶梯：每经过一个素数，函数值跳升 $1$；
在相邻两个素数之间，函数值保持不变。

\subsection{数值表}

下表给出若干处 $\pi(x)$ 的精确值：

\begin{center}
\begin{tabular}{cccccccccc}
\toprule
$x$ & 2 & 3 & 5 & 7 & 10 & 20 & 50 & 100 & 1000 \\
\midrule
$\pi(x)$ & 1 & 2 & 3 & 4 & 4 & 8 & 15 & 25 & 168 \\
\bottomrule
\end{tabular}
\end{center}

\subsection{图形示意}

下图显示 $1\le x\le 30$ 上的 $\pi(x)$：红点为素数处跳跃后的函数值（每点抬升 $1$）。

\begin{center}
\begin{tikzpicture}
\begin{axis}[
    width = 12cm,
    height = 7cm,
    xlabel = {$x$},
    ylabel = {$\pi(x)$},
    title = {素数计数函数 $\pi(x)$（$1 \leq x \leq 30$）},
    xmin=1, xmax=30, ymin=0, ymax=11,
    xtick={2,3,5,7,11,13,17,19,23,29},
    ytick={0,1,2,3,4,5,6,7,8,9,10},
    grid=major,
    grid style={line width=0.3pt, draw=gray!30},
    axis lines=left,
]
\addplot[blue!70!black, thick, const plot, no marks] coordinates {
  (1,0)(2,1)(3,2)(5,3)(7,4)(11,5)(13,6)(17,7)(19,8)(23,9)(29,10)(30,10)
};
\addplot[red, only marks, mark=*, mark size=2pt] coordinates {
  (2,1)(3,2)(5,3)(7,4)(11,5)(13,6)(17,7)(19,8)(23,9)(29,10)
};
\end{axis}
\end{tikzpicture}
\end{center}
% ============================================================
\section{偏移对数积分 \texorpdfstring{$\Li(x)$}{Li(x)}}
\label{sec:offset_logarithmic_integral}
% ============================================================

第~\ref{ch:prime_distribution}~章已引入\textbf{实函数}
\[
  \Li:(1,+\infty)\to\mathbb{R},
  \qquad
  \Li(x)=\int_2^x\frac{\mathrm{d}t}{\log t}
\]
（式~\eqref{eq:Li-def-parts}），并给出分部积分与渐近级数~\eqref{eq:Li-asymptotic}。
本章算术对象以 $\pi,J,\mu,M,\Lambda,\theta,\psi$ 为主；
$\Li$ 仅作 $\pi$ 的实光滑比较主项（定义域 $(1,+\infty)$），细节见第~\ref{ch:logint_Ei}~章。

\noindent\textbf{$\li$ 与 $\Ei$ 不在本章展开。}
Cauchy 主值 $\li(x)$、复幂 $\li(x^{\rho})$、指数积分 $\Ei$，
以及「素数定理用 $\Li$、显式公式用 $\li$、计算用 $\Ei$」的完整论述，
见第~\ref{ch:logint_Ei}~章（置于素数定理提要之后、显式公式之前）。
初等函数章没有复积分与解析延拓工具，不宜同时深讲 $\Li$ 与 $\li$。

\subsection{与 \texorpdfstring{$\pi(x)$}{pi(x)} 的渐近关系}\label{subsec:Li_pi_asymp}

\paragraph{与第~\ref{ch:prime_distribution}~章的分工}
第~\ref{ch:prime_distribution}~章已给出 Legendre / Gauss 经验主项、
$\Li$ 的分部展开、素数定理 $\pi\sim\Li\sim x/\log x$ 的表述（定理~\ref{thm:pnt}）、
以及中等区间对照（表~\ref{tab:pix}、图~\ref{fig:pix_compare}）。
本节只补大尺度数值与变号上界；$\li$ 的严格定义见第~\ref{ch:logint_Ei}~章。

\paragraph{承接：为何 $\Li$ 优于 $x/\log x$}
由第~\ref{ch:prime_distribution}~章~\eqref{eq:Li-parts-1}，
\[
  \Li(x)
  =\frac{x}{\log x}-\frac{2}{\log 2}
  +\int_2^x\frac{\mathrm{d}t}{(\log t)^{2}},
\]
故 $x/\log x$ 是 $\Li$ 的渐近\textbf{首项}；
素数定理断言 $\pi(x)/\Li(x)\to 1$（证明提要见第~\ref{ch:prime_number_theorem}~章）。

\paragraph{$\pi-\Li$ 的符号与首次交叉}
表~\ref{tab:pix} 所示范围内 $\pi-\Li<0$；
对已计算的大范围 $x$（至约 $10^{23}$）亦有 $\pi(x)<\Li(x)$
（文献中常写 $\pi<\li$，因实轴 $\li-\Li$ 为常数，符号与首次交叉问题同类；
见第~\ref{ch:logint_Ei}~章）。
Littlewood（1914）证明 $\pi-\Li$ \textbf{变号无穷多次}。
Skewes（1933，假定 RH）给出上界 $e^{e^{e^{79}}}$（\textbf{Skewes 数}）；
其后可压至约 $1.4\times 10^{316}$（te Riele，1987）。
该上界只保证此前必有一次交叉，并非精确交点。

\paragraph{通向后文}
光滑主项比较用 $\Li$；精确显式公式将用 $\li(x)$、$\li(x^{\rho})$
（第~\ref{ch:logint_Ei}~、\ref{ch:riemann_explicit_formula}~章）。
von Koch 型界与 $M$--RH 预告见本章后文及第~\ref{ch:riemann_explicit_formula}~章。

\subsection{数值比较}
 
\begin{center}
\setlength{\tabcolsep}{4.5pt}
\renewcommand{\arraystretch}{1.4}
\begin{tabular}{r|rrrr}
\toprule
$x$ & $\pi(x)$（精确） & $x/\log x$（Legendre 近似） & $\Li(x)$（Gauss 近似） & 相对误差 $|\pi-\Li|/\pi$ \\
\midrule
$10$ & $4$ & $4.34$ & $5.12$ & $28\%$ \\
$100$ & $25$ & $21.71$ & $29.08$ & $16\%$ \\
$1000$ & $168$ & $144.8$ & $177.6$ & $5.7\%$ \\
$10^6$ & $78498$ & $72382$ & $78628$ & $0.2\%$ \\
\bottomrule
\end{tabular}
\end{center}
 
由表可见：随着 $x$ 增大，$|\pi-\Li|/\pi$ 由小 $x$ 的百分之十几降至 $x=10^6$ 时约 $0.2\%$；
与第~\ref{ch:prime_distribution}~章表~\ref{tab:pix}（至 $10^{10}$，给 $\pi-\lfloor\Li\rfloor$）互相印证，
侧重点不同：彼表强调绝对差与 $\pi<\Li$，本表强调相对误差量级。
此为数值观察，不代替素数定理的证明。

\subsection{图形示意：大尺度渐近比较}
\label{subsec:Li_asymptotic_plot}

中等区间上三者的可分辨比较已见图~\ref{fig:pix_compare}（第~\ref{ch:prime_distribution}~章）。
下图取 $x$ 至约 $10^6$ 量级，展示\textbf{渐近尺度}下的形态。
须区分两条不同的误差：
\begin{itemize}
  \item $\Li(x)-x/\log x$：由 $\Li$ 的渐近展开中\textbf{首项之后的各项}
    （$x/(\log x)^2+2x/(\log x)^3+\cdots$，见~\eqref{eq:Li-parts-1}--\eqref{eq:Li-asymptotic}）
    补上；图中绿线系统低于红线，正是这些高阶项的贡献。
  \item $\pi(x)-\Li(x)$：是素数计数与光滑主项之差，由\textbf{素数定理}保证相对误差
    $|\pi-\Li|/\pi\to 0$（证明提要见第~\ref{ch:prime_number_theorem}~章），
    \textbf{不是}上述级数对 $\Li$ 自身的余项所能解释的；
    其振荡结构由非平凡零点贡献决定，完整表述见第~\ref{ch:riemann_explicit_formula}~章的显式公式。
\end{itemize}
量级上，上表给出 $|\pi-\Li|/\pi$ 由小 $x$ 的百分之十几降至 $x=10^6$ 时约 $0.2\%$；
图中蓝、红线宏观贴合对应的是 $\pi(x)\sim\Li(x)$ 在大尺度下的渐近表现，\textbf{不能}理解为 $\pi\equiv\Li$。

% 非浮动：避免 figure 环境被推迟到第 2 章之外
\begin{center}
\begin{tikzpicture}
\begin{axis}[
  width=0.92\textwidth,
  height=7.2cm,
  xlabel={$x$},
  xmin=0, xmax=1e6,
  ymin=0, ymax=9e4,
  grid=major,
  grid style={line width=0.3pt, draw=gray!30},
  legend pos=north west,
  legend style={font=\small, cells={anchor=west}},
  legend cell align=left,
  tick label style={font=\small},
]
\addplot[blue!75!black, thick, no marks]
  table[x=x, y=pi, col sep=space, header=true] {figures/pi_x_comparison.dat};
\addlegendentry{$\pi(x)$}
\addplot[red!80!black, thick, no marks]
  table[x=x, y=Li, col sep=space, header=true] {figures/pi_x_comparison.dat};
\addlegendentry{$\Li(x)$（Gauss 逼近）}
\addplot[green!55!black, very thick, solid, no marks]
  table[x=x, y=xlog, col sep=space, header=true] {figures/pi_x_comparison.dat};
\addlegendentry{$x/\log x$（Legendre 近似）}
\end{axis}
\end{tikzpicture}
\captionof{figure}{大尺度下 $\pi(x)$、$\Li(x)$ 与 $x/\log x$ 的渐近比较
（横轴约至 $10^6$）。
绿实线为 $x/\log x$（Legendre 近似），系统低于红线 $\Li$：对应 $\Li$ 展开在首项之后的各项；
蓝、红贴合：$|\pi-\Li|/\pi$ 约 $10^{-3}$（上表 $x=10^6$ 为 $0.2\%$）。
中等区间的可分比较见第~\ref{ch:prime_distribution}~章图~\ref{fig:pix_compare}。}
\label{fig:pix_Li_asymptotic}
\end{center}
% ============================================================
\section{Möbius 函数 \texorpdfstring{$\mu(n)$}{mu(n)}}\label{sec:mobius_function}
 
\subsection{定义}
 
Möbius 函数根据正整数的素因子结构，赋予其 $1$、$-1$ 或 $0$ 三种取值。
 
\begin{definition}{}{}
对正整数 $n$，\textbf{Möbius 函数} $\mu(n)$ 定义为：
\[
  \mu(n) \;=\;
  \begin{cases}
    1, & n = 1, \\
    (-1)^k, & n \text{ 是 } k \text{ 个\textbf{不同}素数的乘积}, \\
    0, & n \text{ 含有某个素数的平方作为因子}.
  \end{cases}
\]
\end{definition}
 
三种情况的取值：
\begin{itemize}
  \item 若 $n$ 的素因子各不相同，$\mu(n)$ 根据素因子个数的奇偶性取 $\pm 1$；
  \item 若 $n$ 含平方因子（如 $4 = 2^2$、$12 = 4 \times 3$），则 $\mu(n) = 0$，
        从而在容斥求和中不贡献非零项。
\end{itemize}
$\mu(n)$ 的取值编码了素因子结构在容斥中的符号贡献：
\begin{itemize}
  \item 它是\textbf{积性函数}：$\gcd(m,n)=1$ 时 $\mu(mn)=\mu(m)\mu(n)$，故可逐素数因子处理；
  \item 求和 $\sum_{d\mid n}\mu(d)=[n=1]$ 是容斥原理的算术形式
        （见下节引理~\ref{lem:musum}）。
\end{itemize}
 
\subsection{数值表}
 
\begin{center}
\begin{tabular}{r|rrrrrrrrrrrrr}
\toprule
$n$ & 1 & 2 & 3 & 4 & 5 & 6 & 7 & 8 & 9 & 10 & 11 & 12 & 13 \\
\midrule
$\mu(n)$ & $1$ & $-1$ & $-1$ & $0$ & $-1$ & $1$ & $-1$ & $0$
         & $0$ & $1$ & $-1$ & $0$ & $-1$ \\
\bottomrule
\end{tabular}
\end{center}
 
以 $n=6=2\times3$ 为例：两个不同素因子，$\mu(6)=(-1)^2=1$。
以 $n=4=2^2$ 为例：含平方因子，$\mu(4)=0$。
以 $n=30=2\times3\times5$ 为例：三个不同素因子，$\mu(30)=(-1)^3=-1$。
 
\subsection{Möbius 反演公式}

Möbius 反演给出数论函数与其因子和函数之间的互逆关系。
证明完全初等：先建立 $\mu$ 在因子上的求和性质，再交换双重求和。

\begin{lemma}{}{}\label{lem:musum}
对任意 $n \ge 1$，
\[
\sum_{d \mid n} \mu(d) =
\begin{cases}
1, & n = 1,\\
0, & n > 1.
\end{cases}
\]
\end{lemma}

\begin{proof}
若 $n = 1$，和式仅含 $d = 1$，$\mu(1)=1$，结论成立。
设 $n > 1$，写 $n = p_1^{a_1} \cdots p_k^{a_k}$（$k \ge 1$）。
由 $\mu$ 的积性，
\[
\sum_{d \mid n} \mu(d)
= \prod_{i=1}^{k}
  \Biggl( \sum_{j=0}^{a_i} \mu(p_i^{j}) \Biggr).
\]
而 $\mu(p_i^{0})=1$，$\mu(p_i^{1})=-1$，$\mu(p_i^{j})=0\ (j\ge 2)$，
故每个因子为 $1+(-1)=0$，乘积为 $0$。\qed
\end{proof}

\begin{theorem}{Möbius 反演公式（第一形式）}{}
设 $f, g: \mathbb{N} \to \mathbb{C}$ 为数论函数，且满足
\[
g(n) = \sum_{d \mid n} f(d) \qquad (\forall n \ge 1),
\]
则
\[
f(n)
= \sum_{d \mid n} \mu\!\left(\frac{n}{d}\right) g(d)
= \sum_{d \mid n} \mu(d)\, g\!\left(\frac{n}{d}\right).
\]
\end{theorem}

\begin{proof}
将 $g(d)=\sum_{k\mid d} f(k)$ 代入 $\sum_{d\mid n}\mu(n/d)\,g(d)$：
\[
\sum_{d \mid n} \mu\!\left(\frac{n}{d}\right) g(d)
= \sum_{d \mid n} \mu\!\left(\frac{n}{d}\right) \sum_{k \mid d} f(k).
\]
交换求和次序，令 $d=km$，则 $k\mid n$ 且 $m\mid(n/k)$，得
\[
\sum_{k \mid n} f(k)
\sum_{m \mid \frac{n}{k}} \mu(m).
\]
由引理~\ref{lem:musum}，内层和当且仅当 $n/k=1$（即 $k=n$）时等于 $1$，否则为 $0$。
故整个和式等于 $f(n)$。
右端两个求和写法只是替换 $d\leftrightarrow n/d$，彼此相同。\qed
\end{proof}

\begin{theorem}{Möbius 反演公式（第二形式）}{}
设 $f, g: \mathbb{N} \to \mathbb{C}$ 满足
\[
f(n) = \sum_{d \mid n} \mu(d)\, g\!\left(\frac{n}{d}\right)
\qquad (\forall n \ge 1),
\]
则
\[
g(n) = \sum_{d \mid n} f(d).
\]
\end{theorem}

\begin{proof}
将假设代入 $\sum_{d\mid n} f(d)$：
\[
\sum_{d \mid n} f(d)
= \sum_{d \mid n}
  \sum_{k \mid d}
  \mu(k)\, g\!\left(\frac{d}{k}\right).
\]
令 $d=km$，则 $k\mid n$ 且 $m\mid(n/k)$，交换次序得
\[
\sum_{m \mid n} g(m)
\sum_{k \mid \frac{n}{m}} \mu(k).
\]
再由引理~\ref{lem:musum}，内层和仅当 $n/m=1$（即 $m=n$）时为 $1$，否则为 $0$，
故右端等于 $g(n)$。\qed
\end{proof}

\noindent
第一形式与第二形式互为逆运算：由其一可推其二。

\paragraph{注（Dirichlet 卷积，可选）}
若定义卷积 $(f*g)(n)=\sum_{d\mid n}f(d)g(n/d)$，并记 $1(n)\equiv 1$、
$\varepsilon(1)=1$ 且 $\varepsilon(n)=0\ (n>1)$，
则引理~\ref{lem:musum} 即 $1*\mu=\varepsilon$，
而两种反演合写为
\[
g=1*f \quad\Longleftrightarrow\quad f=\mu*g.
\]
下文用反演时，仅依赖上述初等形式即可。

$\mu$ 作为 Dirichlet 级数系数、与黎曼 $\zeta$ 函数的关系依赖 Euler 乘积，
见第~\ref{ch:riemann_zeta_intro}~章第~\ref{subsec:arith_dirichlet_gen}~节。
本章只用 $\mu$ 的定义与反演。

\subsection{图形示意}

下图用梳状图展示 $\mu(n)$ 的三值结构，与上表一致，并延伸至 $n=30$。
 
\begin{center}
\begin{tikzpicture}
\begin{axis}[
    width = 14cm,
    height = 5cm,
    xlabel = {$n$},
    ylabel = {$\mu(n)$},
    title = {Möbius 函数 $\mu(n)$（$1 \leq n \leq 30$）},
    xmin=0, xmax=31, ymin=-1.3, ymax=1.3,
    xtick={1,2,...,30},
    xticklabel style={font=\tiny},
    ytick={-1,0,1},
    grid=major,
    grid style={line width=0.3pt, draw=gray!30},
    axis x line=center,
    axis y line=left,
    legend style={at={(0.5,-0.28)}, anchor=north,
                  legend columns=3,
                  column sep=1em},
]
\addplot[blue!70!black, only marks, mark=*, mark size=2.2pt] coordinates {
  (1,1)(6,1)(10,1)(14,1)(15,1)(21,1)(22,1)(26,1)
};
\addplot[red!80!black, only marks, mark=square*, mark size=2.2pt] coordinates {
  (2,-1)(3,-1)(5,-1)(7,-1)(11,-1)(13,-1)(17,-1)(19,-1)(23,-1)(29,-1)(30,-1)
};
\addplot[black!60, only marks, mark=*, mark size=2.2pt, fill=green!60] coordinates {
  (4,0)(8,0)(9,0)(12,0)(16,0)(18,0)(20,0)(24,0)(25,0)(27,0)(28,0)
};
\addplot[blue!70!black, thick, ycomb] coordinates {
  (1,1)(6,1)(10,1)(14,1)(15,1)(21,1)(22,1)(26,1)
};
\addplot[red!80!black, thick, ycomb] coordinates {
  (2,-1)(3,-1)(5,-1)(7,-1)(11,-1)(13,-1)(17,-1)(19,-1)(23,-1)(29,-1)(30,-1)
};
\addplot[black!60, thick, ycomb] coordinates {
  (4,0)(8,0)(9,0)(12,0)(16,0)(18,0)(20,0)(24,0)(25,0)(27,0)(28,0)
};
\legend{$\mu(n)=1$, $\mu(n)=-1$, $\mu(n)=0$}
\end{axis}
\end{tikzpicture}
\end{center}

\subsection{Mertens 函数 \texorpdfstring{$M(x)$}{M(x)}：Möbius 函数的累积部分和}\label{subsec:mertens_function}
 
\begin{definition}{}{}
对任意实数 $x \ge 0$，\textbf{Mertens 函数}（Mertens function）
\[
  M(x) \;=\; \sum_{n \le x} \mu(n).
\]
\end{definition}

$M(x)$ 只是 $\mu$ 的部分和：在每个正整数 $n$ 处跳跃量为 $\mu(n)$
（即 $\pm 1$ 或 $0$），在区间 $[n,n+1)$ 上取常值 $M(n)$。
定义完全初等，与 $\zeta$ 无关。

\paragraph{$M(x)$ 函数示意图}
下图将横轴拉长至 $1\le x\le 250$，纵轴关于 $M=0$ 对称取定
（$\lvert M\rvert\le 8$ 在此段内），使正负两侧均匀可见：
$M$ 在 $0$ 上下往返（如 $M(1)=1$、$M(95)=2$、$M(221)=5$，
亦有 $M(199)=-8$），并非恒负。
阶梯含义不变：$\mu(n)=0$ 处不跳，$\mu(n)=\pm 1$ 处升降一格
（可与上一小节 $\mu(n)$ 三值图对照）。

\begin{center}
\begin{tikzpicture}
\begin{axis}[
    width = 14.2cm,
    height = 6.0cm,
    xlabel = {$x$},
    ylabel = {$M(x)$},
    title = {Mertens 函数 $M(x)=\sum_{n\le x}\mu(n)$（$1 \leq x \leq 250$）},
    xmin=0, xmax=255,
    ymin=-9, ymax=9,
    xtick={0,50,100,150,200,250},
    ytick={-8,-6,-4,-2,0,2,4,6,8},
    grid=major,
    grid style={line width=0.3pt, draw=gray!30},
    axis x line=middle,
    axis y line=middle,
    axis line style={gray!60},
    tick label style={font=\small},
    clip=false,
]
% 右连续阶梯：点 $(n,M(n))$；纵轴对称于 $M=0$
\addplot[
    const plot,
    thick,
    blue!70!black,
    mark=*,
    mark size=0.55pt,
    mark options={fill=blue!70!black, draw=blue!70!black},
] coordinates {
  (1,1)(2,0)(3,-1)(4,-1)(5,-2)(6,-1)(7,-2)(8,-2)(9,-2)(10,-1)
  (11,-2)(12,-2)(13,-3)(14,-2)(15,-1)(16,-1)(17,-2)(18,-2)(19,-3)(20,-3)
  (21,-2)(22,-1)(23,-2)(24,-2)(25,-2)(26,-1)(27,-1)(28,-1)(29,-2)(30,-3)
  (31,-4)(32,-4)(33,-3)(34,-2)(35,-1)(36,-1)(37,-2)(38,-1)(39,0)(40,0)
  (41,-1)(42,-2)(43,-3)(44,-3)(45,-3)(46,-2)(47,-3)(48,-3)(49,-3)(50,-3)
  (51,-2)(52,-2)(53,-3)(54,-3)(55,-2)(56,-2)(57,-1)(58,0)(59,-1)(60,-1)
  (61,-2)(62,-1)(63,-1)(64,-1)(65,0)(66,-1)(67,-2)(68,-2)(69,-1)(70,-2)
  (71,-3)(72,-3)(73,-4)(74,-3)(75,-3)(76,-3)(77,-2)(78,-3)(79,-4)(80,-4)
  (81,-4)(82,-3)(83,-4)(84,-4)(85,-3)(86,-2)(87,-1)(88,-1)(89,-2)(90,-2)
  (91,-1)(92,-1)(93,0)(94,1)(95,2)(96,2)(97,1)(98,1)(99,1)(100,1)
  (101,0)(102,-1)(103,-2)(104,-2)(105,-3)(106,-2)(107,-3)(108,-3)(109,-4)(110,-5)
  (111,-4)(112,-4)(113,-5)(114,-6)(115,-5)(116,-5)(117,-5)(118,-4)(119,-3)(120,-3)
  (121,-3)(122,-2)(123,-1)(124,-1)(125,-1)(126,-1)(127,-2)(128,-2)(129,-1)(130,-2)
  (131,-3)(132,-3)(133,-2)(134,-1)(135,-1)(136,-1)(137,-2)(138,-3)(139,-4)(140,-4)
  (141,-3)(142,-2)(143,-1)(144,-1)(145,0)(146,1)(147,1)(148,1)(149,0)(150,0)
  (151,-1)(152,-1)(153,-1)(154,-2)(155,-1)(156,-1)(157,-2)(158,-1)(159,0)(160,0)
  (161,1)(162,1)(163,0)(164,0)(165,-1)(166,0)(167,-1)(168,-1)(169,-1)(170,-2)
  (171,-2)(172,-2)(173,-3)(174,-4)(175,-4)(176,-4)(177,-3)(178,-2)(179,-3)(180,-3)
  (181,-4)(182,-5)(183,-4)(184,-4)(185,-3)(186,-4)(187,-3)(188,-3)(189,-3)(190,-4)
  (191,-5)(192,-5)(193,-6)(194,-5)(195,-6)(196,-6)(197,-7)(198,-7)(199,-8)(200,-8)
  (201,-7)(202,-6)(203,-5)(204,-5)(205,-4)(206,-3)(207,-3)(208,-3)(209,-2)(210,-1)
  (211,-2)(212,-2)(213,-1)(214,0)(215,1)(216,1)(217,2)(218,3)(219,4)(220,4)
  (221,5)(222,4)(223,3)(224,3)(225,3)(226,4)(227,3)(228,3)(229,2)(230,1)
  (231,0)(232,0)(233,-1)(234,-1)(235,0)(236,0)(237,1)(238,0)(239,-1)(240,-1)
  (241,-2)(242,-2)(243,-2)(244,-2)(245,-2)(246,-3)(247,-2)(248,-2)(249,-1)(250,-1)
  (251,-1)
};
\end{axis}
\end{tikzpicture}
\end{center}

本书正文的推导链（$\pi,J,\psi$、显式公式等）并不以 $M$ 为中间步骤。
$M$ 在全书中的特殊角色，是在建立 $\zeta$ 与非平凡零点之后，提供黎曼猜想的一种
\textbf{初等算术等价表述}（陈述如下；完整证明见第~\ref{ch:riemann_explicit_formula}~章第~\ref{sec:mertens_rh}~节）。

后文所称\textbf{黎曼猜想}（RH）是指：经解析延拓后的 $\zeta(s)$，
其全部\textbf{非平凡零点}（落在带 $0<\operatorname{Re}s<1$ 内的零点）
满足 $\operatorname{Re}\rho=\frac12$（临界线）。
记号亦见开篇《符号与约定》。第~\ref{ch:prime_distribution}~章已点出非平凡零点进入显式公式修正项；
RH 本身的解析内容与证明框架留待后文。
在上述含义下，\textbf{后文将证明}有等价关系
\[
  \mathrm{RH}
  \;\Longleftrightarrow\;
  M(x)=O\!\bigl(x^{1/2+\varepsilon}\bigr)
  \quad(\forall\,\varepsilon>0),
\]
即 $\mu$ 的正负号相消充分均匀，使部分和不超过 $x^{1/2+\varepsilon}$ 量级。
这里仅简单介绍 $M$ 在知识谱系中的位置
（完整论证见第~\ref{ch:riemann_explicit_formula}~章第~\ref{sec:mertens-equivalence}~节）。

% ============================================================
\section{黎曼素数幂计数函数 \texorpdfstring{$J(x)$}{J(x)}}\label{sec:riemann_prime_power}
\label{sec:J}

\subsection{定义}
与 $\pi(x)$ 只在素数处计 $1$ 不同，$J(x)$ 对所有素数幂 $p^k$ 加权计数，
在 $p^k$ 处权重为 $1/k$
（黎曼 1859 年将此函数记为 $f$；
定义与后文用法见第~\ref{ch:riemann_explicit_formula}~章第~\ref{sec:integration-by-parts}~节，
原稿记号对照见第~\ref{ch:riemann_original_interpretation}~章）。
按跳跃累加写成定义最为直接：

\begin{definition}{}{}
对实数 $x \ge 2$，\textbf{黎曼素数幂计数函数}定义为
\begin{equation}
  J(x)
  \;=\;
  \sum_{\substack{p^k \le x \\ p\,\text{素数},\, k \ge 1}}
  \frac{1}{k}.
\label{eq:J_def}
\end{equation}
\end{definition}

$J(x)$ 是右连续阶梯函数：仅在素数幂 $p^k$ 处发生跳跃，跳跃量为 $1/k$；
在相邻两个素数幂之间取常值。具体地：
\begin{itemize}
  \item 在素数 $p$（即 $p^1$）处跳跃 $1$；
  \item 在素数平方 $p^2$ 处跳跃 $\frac{1}{2}$；
  \item 在素数立方 $p^3$ 处跳跃 $\frac{1}{3}$；以此类推。
\end{itemize}
注意：$8=2^3$ 的跳跃是 $1/3$，不是 $1/2$。

\subsection{用 \texorpdfstring{$\pi$}{pi} 表示 \texorpdfstring{$J$}{J}}
\label{subsec:J_from_pi}

将 \eqref{eq:J_def} 按幂次 $k$ 分组：固定 $k\ge 1$，满足 $p^k\le x$ 的素数恰为
$p\le x^{1/k}$，共 $\pi(x^{1/k})$ 个，每个贡献 $1/k$。故
\begin{equation}
  J(x)
  \;=\;
  \sum_{k=1}^{N(x)} \frac{1}{k}\,\pi\!\left(x^{1/k}\right)
  \;=\;
  \pi(x) + \frac{1}{2}\pi\!\left(x^{1/2}\right)
         + \frac{1}{3}\pi\!\left(x^{1/3}\right) + \cdots,
\label{eq:J_from_pi}
\end{equation}
其中 $N(x)=\bigl\lfloor\log x/\log 2\bigr\rfloor$。
当 $k>N(x)$ 时 $x^{1/k}<2$，从而 $\pi(x^{1/k})=0$；传统上亦可把上式写到 $\infty$，多出的项恒为零。

式 \eqref{eq:J_from_pi} 与定义 \eqref{eq:J_def} 等价：前者按幂次重排后者。
反过来，从 \eqref{eq:J_from_pi} 读跳跃：点 $x=p^\ell$ 使 $\pi(x^{1/m})$ 发生跳跃，
当且仅当 $p^\ell=q^m$ 对某素数 $q$ 成立，即 $m=\ell$、$q=p$；
于是仅 $m=\ell$ 的项贡献跳跃 $1/\ell$，与定义一致。

\subsection{与 \texorpdfstring{$\pi(x)$}{pi(x)} 的互换关系}\label{subsec:J_pi_mobius}

由 \eqref{eq:J_from_pi} 可用 $\pi$ 表示 $J$；反过来，用 Möbius 反演恢复 $\pi$：
\begin{equation}
\begin{aligned}
  \pi(x)
  &= J(x) - \frac{1}{2}J\!\left(x^{1/2}\right) - \frac{1}{3}J\!\left(x^{1/3}\right)
    - \frac{1}{5}J\!\left(x^{1/5}\right) + \frac{1}{6}J\!\left(x^{1/6}\right) - \cdots \\
  &= \sum_{n=1}^{N(x)}\frac{\mu(n)}{n}\,J\!\left(x^{1/n}\right),
\end{aligned}
\label{eq:pi_from_J}
\end{equation}
其中 $N(x)=\bigl\lfloor\log x/\log 2\bigr\rfloor$；当 $n>N(x)$ 时 $x^{1/n}<2$，从而 $J(x^{1/n})=0$。
含平方因子的 $n$ 有 $\mu(n)=0$，对应项自动消失。

\paragraph{与第~\ref{sec:mobius_function}~节因子和反演的关系}
先分清两种求和（勿与下文混淆）：
\begin{itemize}
  \item \textbf{因子和}（第~\ref{sec:mobius_function}~节）：固定正整数 $n$，对除子求和
    $g(n)=\sum_{d\mid n}f(d)$；Möbius 反演给出 $f(n)=\sum_{d\mid n}\mu(d)\,g(n/d)$。
  \item \textbf{部分和}：如 $M(x)=\sum_{n\le x}\mu(n)$、$\psi(x)=\sum_{n\le x}\Lambda(n)$，
    是对「不超过 $x$ 的自变量」累加，不是对除子求和。
\end{itemize}
式~\eqref{eq:pi_from_J} \textbf{既不是}某对 $f,g$ 的因子和反演的字面代入
（没有「固定 $n$，对 $d\mid n$ 求和」），
\textbf{也不是} $\pi$ 或 $J$ 的部分和公式。

它由 \eqref{eq:J_from_pi} 而来：$J$ 是 $\pi$ 在自变量 $x,x^{1/2},x^{1/3},\ldots$ 上的
线性组合；反演时对同一幂次族加权求和，
计算中仍用第~\ref{sec:mobius_function}~节的恒等式
$\sum_{d\mid m}\mu(d)=[m=1]$。
故与因子和反演\textbf{思想同源、对象不同}：一为除子上的数论函数，一为幂次族 $\{x^{1/n}\}$ 上的阶梯。

\paragraph{推导}
将 \eqref{eq:J_from_pi} 代入 \eqref{eq:pi_from_J} 右端：
\begin{align*}
  \sum_{n\ge 1}\frac{\mu(n)}{n}\,J\!\left(x^{1/n}\right)
  &=
  \sum_{n\ge 1}\frac{\mu(n)}{n}
  \sum_{k\ge 1}\frac{1}{k}\,\pi\!\left(x^{1/(nk)}\right)
  =
  \sum_{n,k\ge 1}\frac{\mu(n)}{nk}\,\pi\!\left(x^{1/(nk)}\right).
\end{align*}
令 $m=nk$（$n$ 跑遍 $m$ 的正除子），则 $\pi(x^{1/m})$ 的系数为
\[
  \frac{1}{m}\sum_{n\mid m}\mu(n)
  =\frac{[m=1]}{m}.
\]
故仅 $m=1$ 的项存活，且系数为 $1$，得左端 $=\pi(x)$，即~\eqref{eq:pi_from_J}。

亦可按跳跃核对：右端在素数 $p$ 处净跳跃为 $1$，在 $p^k$（$k\ge 2$）处净跳跃为 $0$，
与 $\pi(x)$ 一致。

\subsection{图形示意}

下图呈现 $1\le x\le 30$ 上的 $J(x)$：红点为素数处跳跃 $+1$，三角为高次素数幂处跳跃 $+1/k$。

\begin{center}
\begin{tikzpicture}
\begin{axis}[
    width = 13cm,
    height = 7.5cm,
    xlabel = {$x$},
    ylabel = {$J(x)$},
    title = {黎曼素数幂计数函数 $J(x)$（$1 \leq x \leq 30$）},
    xmin=1, xmax=30, ymin=0, ymax=13,
    xtick={2,3,4,5,7,8,9,11,13,16,17,19,23,25,27,29},
    xticklabel style={font=\tiny},
    ytick={0,1,2,3,4,5,6,7,8,9,10,11,12},
    yticklabel style={font=\tiny},
    grid=major,
    grid style={line width=0.3pt, draw=gray!30},
    axis lines=left,
    legend pos=north west,
]
\addplot[blue!70!black, thick, const plot, no marks, forget plot] coordinates {
  (1,0)(2,1.0000)(3,2.0000)(4,2.5000)(5,3.5000)(7,4.5000)
  (8,4.8333)(9,5.3333)(11,6.3333)(13,7.3333)(16,7.5833)
  (17,8.5833)(19,9.5833)(23,10.5833)(25,11.0833)(27,11.4167)
  (29,12.4167)(30,12.4167)
};
\addplot[red, only marks, mark=*, mark size=2pt] coordinates {
  (2,1.0000)(3,2.0000)(5,3.5000)(7,4.5000)(11,6.3333)(13,7.3333)
  (17,8.5833)(19,9.5833)(23,10.5833)(29,12.4167)
};
\addlegendentry{素数 $p$：跳跃 $+1$}
\addplot[pink!80!red, only marks, mark=triangle*, mark size=3pt] coordinates {
  (4,2.5000)(8,4.8333)(9,5.3333)(16,7.5833)(25,11.0833)(27,11.4167)
};
\addlegendentry{素数幂 $p^k$（$k\geq2$）：跳跃 $+1/k$}
\end{axis}
\end{tikzpicture}
\end{center}

与 $\psi(x)$ 的积分关系（$\psi=\int\log t\,\mathrm{d}J$、$J=\int\mathrm{d}\psi/\log t$；
阶梯 Stieltjes 见附录~\ref{app:abel_stieltjes}）
在定义第二 Chebyshev 函数之后给出，见第~\ref{subsec:J_psi_stieltjes}~节。
% ============================================================
\section{von Mangoldt 函数 \texorpdfstring{$\Lambda(n)$}{Lambda(n)}}\label{sec:von_mangoldt}
\label{sec:Lambda}
 
\subsection{定义}
 
\begin{definition}{}{}
对正整数 $n$，\textbf{von Mangoldt 函数} $\Lambda(n)$ 定义为：
\[
  \Lambda(n) \;=\;
  \begin{cases}
    \log p, & n = p^k,\; p \text{ 为素数},\; k \geq 1, \\
    0, & \text{否则}.
  \end{cases}
\]
\end{definition}
$\Lambda(n)$ 只在素数幂处取非零值，且 $p^k$ 处的值与 $p$ 本身相同，均为 $\log p$。
注意 $n=6=2\times3$ 不是素数幂，故 $\Lambda(6)=0$；
$n=4=2^2$ 是素数幂，故 $\Lambda(4)=\log 2$。
其部分和 $\psi(x)=\sum_{n\le x}\Lambda(n)$ 见第~\ref{sec:second_chebyshev}~节；
与 $J$ 的系数关系 $J(x)=\sum_{n\le x}\Lambda(n)/\log n$（$n=1$ 项为 $0$）
在显式公式章展开。

\subsection{数值表}
 
\begin{center}
\begin{tabular}{r|rrrrrrrrrrrrr}
\toprule
$n$ & 1 & 2 & 3 & 4 & 5 & 6 & 7 & 8 & 9 & 10 & 11 & 12 & 13 \\
\midrule
$\Lambda(n)$ & $0$ & $\log 2$ & $\log 3$ & $\log 2$ & $\log 5$ & $0$ & $\log 7$
             & $\log 2$ & $\log 3$ & $0$ & $\log 11$ & $0$ & $\log 13$ \\
\bottomrule
\end{tabular}
\end{center}
 
\subsection{对数恒等式}
任意正整数 $n$ 的自然对数等于其所有正因子上 $\Lambda$ 值之和：
\begin{equation}
  \log n \;=\; \sum_{d \mid n} \Lambda(d).
\label{eq:log_sum_Lambda}
\end{equation}

\begin{proof}
若 $n=1$，两边均为 $0$。设 $n=\prod_{i}p_i^{a_i}$（$a_i\ge 1$）。
则 $\log n=\sum_{i}a_i\log p_i$。
右端 $\sum_{d\mid n}\Lambda(d)$ 仅当 $d$ 为某个 $p_i$ 的正幂时非零：
对固定 $i$，因子 $d=p_i,p_i^2,\ldots,p_i^{a_i}$ 各贡献一次 $\log p_i$，
共 $a_i$ 次。故右端 $=\sum_{i}a_i\log p_i=\log n$。
\end{proof}

例如 $n=12=2^2\times 3$，非零项为 $\Lambda(2)+\Lambda(4)+\Lambda(3)
=\log 2+\log 2+\log 3=\log 12$。

\subsection{von Mangoldt 函数与 Möbius 函数的卷积关系}\label{subsec:Lambda_mu_conv}

恒等式 \eqref{eq:log_sum_Lambda} 恰为第~\ref{sec:mobius_function}~节 Möbius 反演
（第一形式）中的 $g(n)=\log n$、$f(n)=\Lambda(n)$。
由该定理直接得
\begin{equation}
\Lambda(n)
  \;=\;
  \sum_{d\mid n} \mu(d)\,\log\frac{n}{d}.
\label{eq:Lambda_mu_conv}
\end{equation}
再展开 $\log(n/d)=\log n-\log d$：
\[
\sum_{d\mid n}\mu(d)\,\log\frac{n}{d}
  \;=\;
  \log n\sum_{d\mid n}\mu(d)
  \;-\;
  \sum_{d\mid n}\mu(d)\,\log d.
\]
由引理~\ref{lem:musum}，$\sum_{d\mid n}\mu(d)=[n=1]$；
当 $n=1$ 时两边均为 $0$，当 $n>1$ 时 $\log n$ 的系数为 $0$。
故对一切 $n\ge 1$，
\begin{equation}
\Lambda(n)
  \;=\;
  -\sum_{d\mid n} \mu(d)\,\log d.
\label{eq:Lambda_mu_log}
\end{equation}
两式 \eqref{eq:Lambda_mu_conv}、\eqref{eq:Lambda_mu_log} 是 $\Lambda$ 与 $\mu$ 的
Dirichlet 卷积联系，完全在初等数论范围内成立，与 $\zeta$ 无关。

$\Lambda$ 作为 Dirichlet 级数系数、与黎曼 $\zeta$ 的关系依赖 Euler 乘积
（对数导数；\textbf{不}依赖上述 $\mu$ 卷积），
见第~\ref{ch:riemann_zeta_intro}~章第~\ref{subsec:arith_dirichlet_gen}~节。

\subsection{图形示意}

下图在 $1\le n\le 30$ 上显示非零的 $\Lambda(n)$：
蓝钉为素数处 $\log p$，红钉为高次素数幂处（高度仍为 $\log p$，与 $p$ 相同）。

\begin{center}
\begin{tikzpicture}
\begin{axis}[
    width = 13cm,
    height = 7cm,
    xlabel = {$n$},
    ylabel = {$\Lambda(n)$},
    title = {von Mangoldt 函数 $\Lambda(n)$（$1 \leq n \leq 30$）},
    xmin=0, xmax=31, ymin=0, ymax=3.8,
    xtick={1,2,...,30},
    xticklabel style={font=\tiny},
    ytick={0,1,2,3},
    yticklabels={$0$,$1$,$2$,$3$},
    yticklabel style={font=\small},
    grid=major,
    grid style={line width=0.3pt, draw=gray!30},
    axis x line=bottom,
    axis y line=left,
    axis on top,
    legend style={
      at={(0.5,-0.22)},
      anchor=north,
      legend columns=2,
      column sep=1.2em,
      font=\small,
      draw=none,
    },
]
\addplot[blue!70!black, thick, ycomb, mark=*, mark size=2.2pt] coordinates {
  (2,0.6931)(3,1.0986)(5,1.6094)(7,1.9459)
  (11,2.3979)(13,2.5649)(17,2.8332)(19,2.9444)
  (23,3.1355)(29,3.3673)
};
\addlegendentry{素数 $p$：$\Lambda(p)=\log p$}
\addplot[red!80!black, thick, ycomb, mark=*, mark size=2.2pt] coordinates {
  (4,0.6931)(8,0.6931)(9,1.0986)(16,0.6931)(25,1.6094)(27,1.0986)
};
\addlegendentry{素数幂 $p^k$（$k\geq2$）：$\Lambda(p^k)=\log p$}
\end{axis}
\end{tikzpicture}
\end{center}
% ============================================================
\section{第一 Chebyshev 函数 \texorpdfstring{$\theta(x)$}{theta(x)}}\label{sec:first_chebyshev}

本节给出仅对素数取对数权重的 Chebyshev 函数 $\theta$，
以及它与 $\pi$ 的分部求和关系、与素数定理的等价证明提要；与 $\psi$ 的完整恒等式见下节。

\subsection{定义}
 
\begin{definition}{}{}
对实数 $x \geq 0$，\textbf{第一 Chebyshev 函数}定义为不超过 $x$ 的所有素数的对数之和：
\[
  \theta(x) \;=\; \sum_{\substack{p \leq x \\ p\,\text{素数}}} \log p.
\]
\end{definition}
$\theta(x)$ 是右连续阶梯函数，只在素数 $p$ 处跳跃 $\log p$；
在两个相邻素数之间函数值不变。
与 $\psi(x)=\sum_{n\le x}\Lambda(n)$（下节定义）相比，$\theta$ 只在素数处跳跃，
不含高次素数幂 $p^k$（$k\ge 2$）的贡献。

\subsection{数值表}
 
下表列出各素数处 $\theta(x)$ 的精确跳跃后值
（$\theta(x)$ 在相邻素数间保持不变）。\footnote{表中 $e^{\theta(p)}$ 等于不超过 $p$ 的全体素数之积（primorial）。}
 
\begin{center}
\setlength{\tabcolsep}{3.2pt}
\begin{tabular}{r|llllllllll}
\toprule
$p$ & $2$ & $3$ & $5$ & $7$ & $11$ & $13$ & $17$ & $19$ & $23$ & $29$ \\
\midrule
$\theta(p)$ & $0.693$ & $1.792$ & $3.401$ & $5.347$ & $7.745$ & $10.310$
            & $13.143$ & $16.088$ & $19.223$ & $22.590$ \\
$e^{\theta(p)}$ & $2$ & $6$ & $30$ & $210$ & $2310$ & $30030$
            & $510510$ & $9699690$ & $223092870$ & $6469693230$ \\
\bottomrule
\end{tabular}
\end{center}
 
\subsection{与 \texorpdfstring{$\pi(x)$}{pi(x)} 的分部求和关系}\label{subsec:theta_pi_abel}

分部求和（亦称 Abel 求和）给出 $\theta$ 与 $\pi$ 的积分关系。
对右连续阶梯函数，积分 $\int f\,\mathrm{d}A$ 约定为在跳跃点取 $f$ 乘跳跃量之和
（工作定义与系统推演见附录~\ref{app:abel_stieltjes}）。
$\pi$ 在素数 $p$ 处跳跃 $1$，故
$\theta(x)=\sum_{p\le x}\log p=\int_{2^-}^{x}\log t\,\mathrm{d}\pi(t)$；
对右端作分部（附录~\ref{app:abel_stieltjes} 公式~\eqref{eq:app-stieltjes-parts}）得
\begin{equation}
  \theta(x)
  \;=\;
  \pi(x)\log x
  \;-\;
  \int_2^x \frac{\pi(t)}{t}\,\mathrm{d}t.
\label{eq:theta_pi_abel}
\end{equation}


\subsection{与素数定理的等价（证明提要）}\label{subsec:theta_psi_PNT}

\begin{theorem}{}{}
以下两命题等价：
\[
  \pi(x) \sim \frac{x}{\log x}, \qquad
  \theta(x) \sim x.
\]
进而，若记下节的第二 Chebyshev 函数
$\psi(x)=\sum_{n\le x}\Lambda(n)$，则还有
\[
  \theta(x)\sim x
  \quad\Longleftrightarrow\quad
  \psi(x)\sim x,
\]
从而
\[
  \pi(x)\sim\frac{x}{\log x},
  \quad
  \theta(x)\sim x,
  \quad
  \psi(x)\sim x
\]
三者彼此等价。
\end{theorem}

证明提要如下。由 \eqref{eq:theta_pi_abel}，若 $\pi(x)\sim x/\log x$，则主项
$\pi(x)\log x\sim x$，积分项
$\int_2^x\pi(t)/t\,\mathrm{d}t=o(x)$（或更细的 $O(x/\log x)$ 级），
故 $\theta(x)\sim x$。
反过来，若 $\theta(x)\sim x$，则由
$\pi(x)=\int_{2^-}^{x}\mathrm{d}\theta(t)/\log t$
（或对 \eqref{eq:theta_pi_abel} 作标准反演估计）得
$\pi(x)\sim x/\log x$。
细节估计见第~\ref{ch:prime_distribution}~章与附录中素数定理的讨论。

与 $\psi$ 的精确关系（公式先写，定义与证明见下节）：按素数幂分组有
\[
  \psi(x)
  \;=\;
  \sum_{k\ge 1}\theta\!\left(x^{1/k}\right)
  \;=\;
  \theta(x)+\theta\!\left(x^{1/2}\right)+\theta\!\left(x^{1/3}\right)+\cdots,
\]
且 $\psi(x)=\theta(x)+O(\sqrt{x}\,\log x)$。
因此 $\theta(x)-x=o(x)$ 当且仅当 $\psi(x)-x=o(x)$，即
$\theta\sim x\iff\psi\sim x$。
编号公式与证明见第~\ref{sec:second_chebyshev}~节 \eqref{eq:psi_sum_theta}、\eqref{eq:psi_theta_asymp}。

\subsection{Chebyshev 的初等估计}
 
在黎曼猜想出现之前，Chebyshev（1852）已用初等方法证明了：
存在常数 $A < 1 < B$ 使得对充分大的 $x$，
\[
A x \;<\; \theta(x) \;<\; B x.
\]
例如可取 $A = 0.921$，$B = 1.106$。
由 \eqref{eq:theta_pi_abel}，这给出
$\pi(x)=O(x/\log x)$ 与 $\pi(x)=\Omega(x/\log x)$，
是素数定理 $\pi\sim x/\log x$ 的第一个实质性逼近（尚未达到 $\sim$）。
 
\subsection{图形示意}

下图描绘 $1\le x\le 30$ 上的 $\theta(x)$：红点为素数处跳跃后的函数值。

\begin{center}
\begin{tikzpicture}
\begin{axis}[
    width = 13cm,
    height = 7.5cm,
    xlabel = {$x$},
    ylabel = {$\theta(x)$},
    title = {第一 Chebyshev 函数 $\theta(x)$（$1 \leq x \leq 30$）},
    xmin=1, xmax=30, ymin=0, ymax=26,
    xtick={2,3,5,7,11,13,17,19,23,29},
    ytick={0,5,10,15,20,25},
    yticklabels={$0$,$5$,$10$,$15$,$20$,$25$},
    yticklabel style={font=\small},
    grid=major,
    grid style={line width=0.3pt, draw=gray!30},
    axis lines=left,
    legend style={
      at={(0.02,0.98)},
      anchor=north west,
      font=\small,
      draw=none,
      fill=white,
      fill opacity=0.85,
      text opacity=1,
    },
]
\addplot[blue!70!black, thick, const plot, no marks] coordinates {
  (1,0.0)(2,0.6931)(3,1.7918)(5,3.4012)(7,5.3471)
  (11,7.745)(13,10.31)(17,13.1432)(19,16.0876)(23,19.2231)(29,22.5904)(30,22.5904)
};
\addlegendentry{$\theta(x)$（阶梯）}
\addplot[red!80!black, only marks, mark=*, mark size=2.5pt] coordinates {
  (2,0.6931)(3,1.7918)(5,3.4012)(7,5.3471)
  (11,7.745)(13,10.31)(17,13.1432)(19,16.0876)(23,19.2231)(29,22.5904)
};
\addlegendentry{素数跳跃点}
\end{axis}
\end{tikzpicture}
\end{center}
% ============================================================
\section{第二 Chebyshev 函数 \texorpdfstring{$\psi(x)$}{psi(x)}}\label{sec:second_chebyshev}

第二 Chebyshev 函数 $\psi(x)$ 是本章最便于与 $-\zeta'/\zeta$ 对接的阶梯对象：
它是 $\Lambda$ 的部分和，后文 Perron 积分与显式公式多以 $\psi$ 为主角。
本节给出定义、与 $\theta$、$J$ 的初等恒等式，以及与 $\theta$ 的对比图；
解析侧（生成函数、Perron、$\psi\sim x$、显式公式）只作前指。

\subsection{定义}
 
\begin{definition}{}{}
对实数 $x \geq 2$，\textbf{第二 Chebyshev 函数}定义为
\begin{equation}
  \psi(x)
  \;=\;
  \sum_{\substack{p^k \leq x \\ p\,\text{素数},\, k \geq 1}} \log p
  \;=\;
  \sum_{n \leq x} \Lambda(n).
\label{eq:psi_def}
\end{equation}
\end{definition}

两式同一：由第~\ref{sec:Lambda}~节，$\Lambda(n)$ 恰在 $n=p^k$ 时取值 $\log p$，
其余为 $0$，故对 $n\le x$ 求和即得按素数幂的写法。

$\psi$ 是右连续阶梯函数：仅在素数幂 $p^k$ 处发生跳跃，跳跃量为 $\log p$
（注意：与 $k$ 无关，例如 $4=2^2$、$8=2^3$ 均跳 $\log 2$）。
在相邻两个素数幂之间取常值。
与 $\theta(x)$（仅在素数 $p$ 处跳 $\log p$）相比，$\psi$ 在 $p^2,p^3,\ldots$ 处
仍有额外跳跃，故恒有 $\psi(x)\ge\theta(x)$。

\subsection{数值表}

下表按跳跃点列出 $\psi$ 的累加过程（$\approx$ 为三位小数；精确值见图中坐标）。

\begin{center}
\renewcommand{\arraystretch}{1.3}
\begin{tabular}{r|lll}
\toprule
跳跃点 & 类型 & 跳跃量 & $\psi$ 的新值 \\
\midrule
$2$ & 素数 $2^1$ & $\log 2 \approx 0.693$ & $0.693$ \\
$3$ & 素数 $3^1$ & $\log 3 \approx 1.099$ & $1.792$ \\
$4$ & 素数幂 $2^2$ & $\log 2 \approx 0.693$ & $2.485$ \\
$5$ & 素数 $5^1$ & $\log 5 \approx 1.609$ & $4.094$ \\
$7$ & 素数 $7^1$ & $\log 7 \approx 1.946$ & $6.040$ \\
$8$ & 素数幂 $2^3$ & $\log 2 \approx 0.693$ & $6.733$ \\
$9$ & 素数幂 $3^2$ & $\log 3 \approx 1.099$ & $7.832$ \\
$11$ & 素数 $11^1$ & $\log 11 \approx 2.398$ & $10.230$ \\
$13$ & 素数 $13^1$ & $\log 13 \approx 2.565$ & $12.795$ \\
$16$ & 素数幂 $2^4$ & $\log 2 \approx 0.693$ & $13.488$ \\
$17$ & 素数 $17^1$ & $\log 17 \approx 2.833$ & $16.321$ \\
$19$ & 素数 $19^1$ & $\log 19 \approx 2.944$ & $19.266$ \\
$23$ & 素数 $23^1$ & $\log 23 \approx 3.135$ & $22.401$ \\
$25$ & 素数幂 $5^2$ & $\log 5 \approx 1.609$ & $24.010$ \\
$27$ & 素数幂 $3^3$ & $\log 3 \approx 1.099$ & $25.109$ \\
$29$ & 素数 $29^1$ & $\log 29 \approx 3.367$ & $28.476$ \\
\bottomrule
\end{tabular}
\end{center}

\subsection{与 \texorpdfstring{$\theta(x)$}{theta(x)} 的关系}\label{subsec:psi_theta_rel}

上节已写出 \eqref{eq:psi_sum_theta}、\eqref{eq:psi_theta_asymp}；此处给出证明与反演。

\begin{theorem}{}{}
对 $x \ge 1$，
\begin{equation}
  \psi(x)
  \;=\;
  \sum_{k\ge 1}\theta\!\left(x^{1/k}\right)
  \;=\;
  \theta(x)+\theta\!\left(x^{1/2}\right)+\theta\!\left(x^{1/3}\right)+\cdots,
\label{eq:psi_sum_theta}
\end{equation}
其中当 $x^{1/k}<2$ 时 $\theta(x^{1/k})=0$，故求和实际上有限
（至多到 $k=\lfloor\log x/\log 2\rfloor$）。
\end{theorem}

\begin{proof}
由定义 \eqref{eq:psi_def}，
\[
\psi(x)
  \;=\;
  \sum_{k\ge 1}\sum_{\substack{p \\ p^k\le x}}\log p
  \;=\;
  \sum_{k\ge 1}\sum_{p\le x^{1/k}}\log p
  \;=\;
  \sum_{k\ge 1}\theta\!\left(x^{1/k}\right).
\]
\end{proof}

在平凡界 $\theta(y)=O(y\log y)$ 下，$\theta(x^{1/2})=O(\sqrt{x}\,\log x)$，
$k\ge 3$ 的项更小，故
\begin{equation}
  \psi(x)
  \;=\;
  \theta(x)
  \;+\;
  O\!\left(\sqrt{x}\,\log x\right).
\label{eq:psi_theta_asymp}
\end{equation}
从而 $\theta(x)\sim x$ 当且仅当 $\psi(x)\sim x$（见第~\ref{subsec:theta_psi_PNT}~节证明提要）。

反过来，由 \eqref{eq:psi_sum_theta} 作 Möbius 反演（与第~\ref{subsec:J_pi_mobius}~节
$J\leftrightarrow\pi$ 同一形态：变量 $x\mapsto x^{1/k}$），得
\begin{equation}
  \theta(x)
  \;=\;
  \sum_{k\ge 1}\mu(k)\,\psi\!\left(x^{1/k}\right).
\label{eq:theta_from_psi}
\end{equation}
证明轮廓：将 \eqref{eq:psi_sum_theta} 代入右端，
\[
\sum_{k\ge 1}\mu(k)\,\psi\!\left(x^{1/k}\right)
  \;=\;
  \sum_{k,m\ge 1}\mu(k)\,\theta\!\left(x^{1/(km)}\right);
\]
令 $\ell=km$，则 $\theta(x^{1/\ell})$ 的系数为
$\sum_{k\mid\ell}\mu(k)=[\ell=1]$，故仅 $\ell=1$ 存活，得 $\theta(x)$。

\subsection{\texorpdfstring{$J(x)$}{J(x)} 与 \texorpdfstring{$\psi(x)$}{psi(x)} 的积分关系}\label{subsec:J_psi_stieltjes}

第~\ref{sec:J}~节已定义 $J(x)=\sum_{p^k\le x}1/k$（在 $p^k$ 处跳跃 $1/k$）；
本节的 $\psi(x)$ 在同一批点处跳跃 $\log p$。
对右连续阶梯函数，积分 $\int f\,\mathrm{d}A$ 表示跳跃处 $f$ 与跳跃量的乘积和
（亦称阶梯上的 Stieltjes 积分；定义与分部公式见附录~\ref{app:abel_stieltjes}）。
二者由此互逆，全部在初等范围内成立，不依赖 $\zeta(s)$。

\begin{theorem}{}{}
对任意 $x > 1$，
\begin{equation}
\psi(x) \;=\; \int_{2^-}^{x} \log t \; \mathrm{d}J(t),
\label{eq:psi_from_J_stieltjes}
\end{equation}
\begin{equation}
J(x) \;=\; \int_{2^-}^{x} \frac{1}{\log t} \; \mathrm{d}\psi(t),
\label{eq:J_from_psi_stieltjes}
\end{equation}
\end{theorem}

\begin{proof}
由附录~\ref{app:abel_stieltjes} 定义~\ref{def:app-stieltjes}，积分化为跳跃有限和。
对 \eqref{eq:psi_from_J_stieltjes}：在 $t=p^k$ 处 $J$ 的跳跃为 $1/k$，故
\[
\int_{2^-}^{x} \log t \; \mathrm{d}J(t)
  \;=\; \sum_{p^k \le x} \log(p^k) \cdot \frac{1}{k}
  \;=\; \sum_{p^k \le x} \log p
  \;=\; \psi(x).
\]
对 \eqref{eq:J_from_psi_stieltjes}：在 $t=p^k$ 处 $\psi$ 的跳跃为 $\log p$，故
\[
\int_{2^-}^{x} \frac{1}{\log t} \; \mathrm{d}\psi(t)
  \;=\; \sum_{p^k \le x} \frac{\log p}{\log(p^k)}
  \;=\; \sum_{p^k \le x} \frac{\log p}{k\log p}
  \;=\; \sum_{p^k \le x} \frac{1}{k}
  \;=\; J(x).
\]
\end{proof}

由附录~\ref{app:abel_stieltjes} 分部公式~\eqref{eq:app-stieltjes-parts}，
\eqref{eq:psi_from_J_stieltjes} 还有不显式出现 $\mathrm{d}J$ 的写法：
\begin{equation}
\psi(x) \;=\; \log x \cdot J(x) \;-\; \int_1^x \frac{J(t)}{t} \, \mathrm{d}t.
\label{eq:psi_J_parts}
\end{equation}
从 $\psi$ 表出 $J$ 的分部形式
（$J(x)=\psi(x)/\log x+\int_2^x\psi(t)/(t(\log t)^2)\,\mathrm{d}t$）
见第~\ref{ch:riemann_explicit_formula}~章第~\ref{sec:integration-by-parts}~节
及附录~\ref{app:abel_stieltjes} 式~\eqref{eq:app-J-from-psi-parts}。

跳跃权重对照：$J$ 在 $p^k$ 处为 $1/k$，$\psi$ 为 $\log p$；
因子 $\log t$ 与 $1/\log t$ 恰好互逆，故 \eqref{eq:psi_from_J_stieltjes}、\eqref{eq:J_from_psi_stieltjes} 互为逆变换。
另有系数写法
\[
J(x)=\sum_{2\le n\le x}\frac{\Lambda(n)}{\log n},
\]
与 \eqref{eq:J_from_psi_stieltjes} 在跳跃处逐项一致（$n=p^k$ 时 $\Lambda(n)/\log n=1/k$）。

\subsection{与后文解析对象的衔接}

以下关系\textbf{不在本章证明}，仅标明 $\psi$ 在全书主链上的位置：

\begin{itemize}
  \item \textbf{生成函数}（$\operatorname{Re}s>1$）：
    $\sum_{n=1}^{\infty}\Lambda(n)n^{-s}=-\zeta'(s)/\zeta(s)$，
    故 $\psi(x)=\sum_{n\le x}\Lambda(n)$ 是该级数的部分和
    （第~\ref{ch:riemann_zeta_intro}~章第~\ref{subsec:arith_dirichlet_gen}~节）。
  \item \textbf{Perron 积分}：
    $\psi(x)=\frac{1}{2\pi i}\int_{c-i\infty}^{c+i\infty}
    \bigl(-\zeta'(s)/\zeta(s)\bigr)\,x^s/s\,\mathrm{d}s$（$c>1$）
    （第~\ref{ch:integral_transforms}~章）。
  \item \textbf{素数定理}：$\psi(x)\sim x$ 与 $\pi(x)\sim x/\log x$、$\theta(x)\sim x$ 等价
    （第~\ref{subsec:theta_psi_PNT}~节证明提要；证明见第~\ref{ch:prime_number_theorem}~章）。
  \item \textbf{显式公式}：围道左移后得 $\psi_0(x)$ 的零点展开
    （跳跃点取左右平均；见《符号与约定》中 $\psi_0$ 与第~\ref{ch:riemann_explicit_formula}~章）。
\end{itemize}

\subsection{图形示意}

下图在 $1\le x\le 30$ 上对比 $\psi$（蓝）与 $\theta$（红）：
二者在素数处同步抬升；蓝色额外在 $4,8,9,16,25,27$ 等 $p^k$（$k\ge 2$）处再跳 $\log p$
（三角标记），故蓝阶始终不低于红阶。

\begin{center}
\begin{tikzpicture}
\begin{axis}[
    width = 13cm,
    height = 7.5cm,
    xlabel = {$x$},
    ylabel = {},
    title = {$\psi(x)$（蓝）与 $\theta(x)$（红）对比（$1 \leq x \leq 30$）},
    xmin=1, xmax=30, ymin=0, ymax=31,
    xtick={2,3,4,5,7,8,9,11,13,16,17,19,23,25,27,29},
    xticklabel style={font=\tiny},
    ytick={0,5,10,15,20,25,30},
    grid=major,
    grid style={line width=0.3pt, draw=gray!30},
    axis lines=left,
    legend style={
      at={(0.02,0.98)},
      anchor=north west,
      font=\small,
      draw=black!40,
      fill=white,
      fill opacity=0.92,
      text opacity=1,
      row sep=2pt,
    },
    legend cell align=left,
]
\addplot[blue!70!black, thick, const plot, no marks] coordinates {
  (1,0.0)(2,0.6931)(3,1.7918)(4,2.4849)(5,4.0943)
  (7,6.0403)(8,6.7334)(9,7.832)(11,10.2299)(13,12.7949)
  (16,13.488)(17,16.3212)(19,19.2657)(23,22.4012)
  (25,24.0106)(27,25.1092)(29,28.4765)(30,28.4765)
};
\addlegendentry{$\psi(x)$（含素数幂）}
\addplot[forget plot, blue!70!black, only marks, mark=*, mark size=2.5pt] coordinates {
  (2,0.6931)(3,1.7918)(5,4.0943)(7,6.0403)
  (11,10.2299)(13,12.7949)(17,16.3212)(19,19.2657)(23,22.4012)(29,28.4765)
};
\addplot[cyan!60!black, only marks, mark=triangle*, mark size=3pt] coordinates {
  (4,2.4849)(8,6.7334)(9,7.832)(16,13.488)(25,24.0106)(27,25.1092)
};
\addlegendentry{$p^k$（$k\ge 2$）额外跳跃}
\addplot[red!75!black, thick, const plot, no marks] coordinates {
  (1,0.0)(2,0.6931)(3,1.7918)(5,3.4012)(7,5.3471)
  (11,7.745)(13,10.31)(17,13.1432)(19,16.0876)(23,19.2231)(29,22.5904)(30,22.5904)
};
\addlegendentry{$\theta(x)$（仅素数）}
\addplot[forget plot, red!75!black, only marks, mark=*, mark size=2pt] coordinates {
  (2,0.6931)(3,1.7918)(5,3.4012)(7,5.3471)
  (11,7.745)(13,10.31)(17,13.1432)(19,16.0876)(23,19.2231)(29,22.5904)
};
% 显式保证红线进入图例（仅 mark 的图用 forget plot，避免挤占图例项）
\end{axis}
\end{tikzpicture}
\end{center}
% ============================================================
\section{各函数的关系}\label{sec:relations_seven_functions}
 
\subsection{分类汇总}
 
\begin{center}
\renewcommand{\arraystretch}{1.6}
\begin{tabular}{llllll}
\toprule
函数 & 类别 & 定义域 & 跳跃位置 & 跳跃量 & 图像特征 \\
\midrule
$\pi(x)$ & 素数计数 & $\R_{\geq 0}$ & 素数 $p$ & $1$ & 整数阶梯 \\
$J(x)$ & 素数计数 & $\R_{>0}$ & 素数幂 $p^k$ & $1/k$ & 有理阶梯 \\
$\theta(x)$ & 素数计数 & $\R_{\geq 0}$ & 素数 $p$ & $\log p$ & 对数阶梯 \\
$\psi(x)$ & 素数计数 & $\R_{\geq 0}$ & 素数幂 $p^k$ & $\log p$ & 对数阶梯 \\
$\mu(n)$ & 积性算术函数 & $\N$ & 每个正整数 & $\in\{-1,0,1\}$ & 双向跳动 \\
$M(x)$ & $\mu$ 的部分和 & $\R_{\geq 0}$ & 正整数处 & $\mu(n)$ & 游走型阶梯 \\
$\Lambda(n)$ & 算术函数（\textbf{非积性}） & $\N$ & 素数幂（其余为 $0$） & $\log p$ & 单向脉冲 \\
$\Li(x)$ & 解析近似 & $\R_{>1}$ & 无（连续） & — & 光滑曲线 \\
\bottomrule
\end{tabular}
\end{center}
 
\newpage
\subsection{初等数论函数间的主要数学关系}
\label{sec:elementary_spectrum_map}

本章已建立素数计数函数 $\pi,J,\theta,\psi$、算术函数 $\mu$、$M$、$\Lambda$ 以及主项 $\Li(x)$ 之间的初等关系。
下图采用与第~\ref{ch:spectrum_map}~章全文图~\ref{fig:function-relations}
\textbf{同一底图}（与第~\ref{ch:spectrum_map}~章全文谱系图一致）：
\textbf{彩色}标出截至本章已建立的主要对象与关系；\textbf{灰白}标出后文才讨论的对象与关系
（节点、连线及其上公式均保留，仅作灰白化，不删除）。
读图时宜将彩色元素与本章相应小节对照，把图上的箭头读成已证公式的索引。

% 图独占一页；题注文字排到下一页，避免与图例挤在同页
% （同一底图三次着色之第一次；\spectrumStage=2 为内部进度码，非章号）
\clearpage
\begin{figure}[p]
\centering
\def\spectrumStage{2}
\resizebox{\textwidth}{!}{%
% 黎曼猜想知识谱系关系图（A-2 现行底图 + 分阶段灰白）
% 用法：\def\spectrumStage{2|8|16}% 黎曼猜想知识谱系关系图（A-2 现行底图 + 分阶段灰白）
% 用法：\def\spectrumStage{2|8|16}% 黎曼猜想知识谱系关系图（A-2 现行底图 + 分阶段灰白）
% 用法：\def\spectrumStage{2|8|16}\input{图谱/谱系关系图-core.tex}
% 几何与 黎曼猜想知识谱系关系图-A-2.tex 同步；未涉及者仅灰白化（含图例）
%
% \spectrumStage 仅为内部进度码（历史命名），不是章号：
%   2  → 第一次着色：第2章 初等算术层（fig:spectrum-stage2）
%   8  → 第二次着色：第9章 延拓/函数方程（fig:spectrum-stage8）
%  16  → 第三次着色：第15章 全文总图（fig:function-relations）
% 正文表述用「第一次 / 第二次 / 第三次着色」或章标签，勿写「第8章图」「第16章图」。
\providecommand{\spectrumStage}{16}
\makeatletter
\colorlet{cDim}{gray!16}\colorlet{dDim}{gray!48}\colorlet{tDim}{gray!52}
\ifnum\spectrumStage=2
  \colorlet{cPrime}{blue!12}\colorlet{dPrime}{black}\colorlet{tPrime}{black}
  \colorlet{cMu}{teal!15}\colorlet{dMu}{teal!60!black}\colorlet{tMu}{black}
  \colorlet{cLi}{yellow!25}\colorlet{dLi}{orange!70!black}\colorlet{tLi}{black}
  \colorlet{cGammaU}{gray!16}\colorlet{cGammaL}{gray!16}\colorlet{dGamma}{gray!48}\colorlet{tGamma}{gray!52}
  \colorlet{cZetaU}{gray!16}\colorlet{cZetaL}{gray!16}\colorlet{tZeta}{gray!52}
  \colorlet{cXi}{gray!16}\colorlet{dXi}{gray!48}\colorlet{tXi}{gray!52}
  \colorlet{cNT}{gray!16}\colorlet{dNT}{gray!48}\colorlet{tNT}{gray!52}
  \colorlet{cProd}{gray!16}\colorlet{dProd}{gray!48}\colorlet{tProd}{gray!52}
  \colorlet{cRH}{gray!16}\colorlet{dRH}{gray!48}\colorlet{tRH}{gray!52}
  \colorlet{eArith}{blue!70!black}\colorlet{eElem}{violet!80!black}\colorlet{eDef}{black!80}
  \colorlet{ePNT}{brown!70!black}\colorlet{eAnal}{gray!42}\colorlet{eExpl}{gray!42}
  \colorlet{eExt}{gray!42}\colorlet{eZero}{gray!42}\colorlet{eLate}{gray!42}
  \colorlet{lgPrime}{blue!12}\colorlet{lgMu}{teal!15}\colorlet{lgLi}{yellow!25}
  \colorlet{lgXi}{gray!22}\colorlet{lgRH}{gray!22}\colorlet{lgGU}{gray!22}\colorlet{lgGL}{gray!22}
  \colorlet{lgEArith}{blue!70!black}\colorlet{lgEDef}{black!80}\colorlet{lgEElem}{violet!80!black}
  \colorlet{lgEZero}{gray!45}\colorlet{lgEAnal}{gray!45}\colorlet{lgEPNT}{brown!70!black}\colorlet{lgEExt}{gray!45}
  \colorlet{lgTprime}{black}\colorlet{lgTmu}{black}\colorlet{lgTli}{black}
  \colorlet{lgTxi}{gray!50}\colorlet{lgTrh}{gray!50}\colorlet{lgTgamma}{gray!50}
  \colorlet{lgTarith}{black}\colorlet{lgTdef}{black}\colorlet{lgTelem}{black}
  \colorlet{lgTzero}{gray!50}\colorlet{lgTanal}{gray!50}\colorlet{lgTpnt}{black}\colorlet{lgText}{gray!50}
\else
\ifnum\spectrumStage=8
  \colorlet{cPrime}{blue!12}\colorlet{dPrime}{black}\colorlet{tPrime}{black}
  \colorlet{cMu}{teal!15}\colorlet{dMu}{teal!60!black}\colorlet{tMu}{black}
  \colorlet{cLi}{yellow!25}\colorlet{dLi}{orange!70!black}\colorlet{tLi}{black}
  \colorlet{cGammaU}{blue!15}\colorlet{cGammaL}{red!15}\colorlet{dGamma}{black}\colorlet{tGamma}{black}
  \colorlet{cZetaU}{blue!15}\colorlet{cZetaL}{red!15}\colorlet{tZeta}{black}
  \colorlet{cXi}{gray!16}\colorlet{dXi}{gray!48}\colorlet{tXi}{gray!52}
  \colorlet{cNT}{gray!16}\colorlet{dNT}{gray!48}\colorlet{tNT}{gray!52}
  \colorlet{cProd}{gray!16}\colorlet{dProd}{gray!48}\colorlet{tProd}{gray!52}
  \colorlet{cRH}{gray!16}\colorlet{dRH}{gray!48}\colorlet{tRH}{gray!52}
  \colorlet{eArith}{blue!70!black}\colorlet{eElem}{violet!80!black}\colorlet{eDef}{black!80}
  \colorlet{ePNT}{brown!70!black}\colorlet{eAnal}{green!50!black}\colorlet{eExpl}{gray!42}
  \colorlet{eExt}{red!70!black}\colorlet{eZero}{gray!42}\colorlet{eLate}{gray!42}
  \colorlet{lgPrime}{blue!12}\colorlet{lgMu}{teal!15}\colorlet{lgLi}{yellow!25}
  \colorlet{lgXi}{gray!22}\colorlet{lgRH}{gray!22}\colorlet{lgGU}{blue!15}\colorlet{lgGL}{red!15}
  \colorlet{lgEArith}{blue!70!black}\colorlet{lgEDef}{black!80}\colorlet{lgEElem}{violet!80!black}
  \colorlet{lgEZero}{gray!45}\colorlet{lgEAnal}{green!50!black}\colorlet{lgEPNT}{brown!70!black}\colorlet{lgEExt}{red!70!black}
  \colorlet{lgTprime}{black}\colorlet{lgTmu}{black}\colorlet{lgTli}{black}
  \colorlet{lgTxi}{gray!50}\colorlet{lgTrh}{gray!50}\colorlet{lgTgamma}{black}
  \colorlet{lgTarith}{black}\colorlet{lgTdef}{black}\colorlet{lgTelem}{black}
  \colorlet{lgTzero}{gray!50}\colorlet{lgTanal}{black}\colorlet{lgTpnt}{black}\colorlet{lgText}{black}
\else
  \colorlet{cPrime}{blue!12}\colorlet{dPrime}{black}\colorlet{tPrime}{black}
  \colorlet{cMu}{teal!15}\colorlet{dMu}{teal!60!black}\colorlet{tMu}{black}
  \colorlet{cLi}{yellow!25}\colorlet{dLi}{orange!70!black}\colorlet{tLi}{black}
  \colorlet{cGammaU}{blue!15}\colorlet{cGammaL}{red!15}\colorlet{dGamma}{black}\colorlet{tGamma}{black}
  \colorlet{cZetaU}{blue!15}\colorlet{cZetaL}{red!15}\colorlet{tZeta}{black}
  \colorlet{cXi}{cyan!15}\colorlet{dXi}{cyan!60!black}\colorlet{tXi}{black}
  \colorlet{cNT}{cyan!15}\colorlet{dNT}{cyan!60!black}\colorlet{tNT}{black}
  \colorlet{cProd}{cyan!15}\colorlet{dProd}{cyan!60!black}\colorlet{tProd}{black}
  \colorlet{cRH}{orange!15}\colorlet{dRH}{orange!70!black}\colorlet{tRH}{black}
  \colorlet{eArith}{blue!70!black}\colorlet{eElem}{violet!80!black}\colorlet{eDef}{black!80}
  \colorlet{ePNT}{brown!70!black}\colorlet{eAnal}{green!50!black}\colorlet{eExpl}{green!50!black}
  \colorlet{eExt}{red!70!black}\colorlet{eZero}{orange!80!black}\colorlet{eLate}{black!80}
  \colorlet{lgPrime}{blue!12}\colorlet{lgMu}{teal!15}\colorlet{lgLi}{yellow!25}
  \colorlet{lgXi}{cyan!15}\colorlet{lgRH}{orange!15}\colorlet{lgGU}{blue!15}\colorlet{lgGL}{red!15}
  \colorlet{lgEArith}{blue!70!black}\colorlet{lgEDef}{black!80}\colorlet{lgEElem}{violet!80!black}
  \colorlet{lgEZero}{orange!80!black}\colorlet{lgEAnal}{green!50!black}\colorlet{lgEPNT}{brown!70!black}\colorlet{lgEExt}{red!70!black}
  \colorlet{lgTprime}{black}\colorlet{lgTmu}{black}\colorlet{lgTli}{black}
  \colorlet{lgTxi}{black}\colorlet{lgTrh}{black}\colorlet{lgTgamma}{black}
  \colorlet{lgTarith}{black}\colorlet{lgTdef}{black}\colorlet{lgTelem}{black}
  \colorlet{lgTzero}{black}\colorlet{lgTanal}{black}\colorlet{lgTpnt}{black}\colorlet{lgText}{black}
\fi\fi
\makeatother

\begin{tikzpicture}[scale=1.0, transform shape,
  box/.style    = {draw, rounded corners=5pt, minimum width=2.8cm,
                   minimum height=0.9cm, align=center, font=\small, thick},
  zbox/.style   = {draw, rounded corners=5pt, minimum width=2.8cm,
                   minimum height=0.9cm, align=center, font=\small,
                   thick, fill=gray!18},
  splitbox/.style = {draw, rounded corners=5pt, minimum width=7cm,
                     minimum height=2.2cm, align=center, font=\small, thick,
                     rectangle split, rectangle split parts=2,
                     rectangle split part fill={red!8, white}},
  arr/.style    = {-{Stealth[length=6pt,width=5pt]}, thick},
  barr/.style   = {{Stealth[length=6pt,width=5pt]}-{Stealth[length=6pt,width=5pt]}, thick},
  % 线型：虚线仅用于 N(T)→RH、ξ↔RH；其余为实线
  proven/.style = {solid},
  cond/.style   = {dashed},
  lab/.style    = {font=\tiny, align=center, inner sep=3pt},
  rhbox/.style  = {draw=dRH, rounded corners=5pt, minimum width=8cm,
                   minimum height=1.0cm, align=center, font=\small,
                   fill=cRH, text=tRH, thick}
]

% 布局与 A-2 同步：左移 1.0cm、上移 2.0cm
% 左右略放宽，避免 Γ/ζ 自环旁公式标注被裁切
\path[use as bounding box]
  (-10.2, -26.2) rectangle (10.0, 7.4);

\begin{scope}[shift={(-1.0,2.0)}]

\node[font=\bfseries\Large, align=center] at (0, 4.35)
  {黎曼猜想知识谱系关系图};

% 主图单独上移 1cm；标题与图例位置不动
\begin{scope}[shift={(-2.75,-1.2)}]
%% ==================== 节点分层 ====================

% 第一层：素数计数函数 (y=0)
\node[box, fill=cPrime, draw=dPrime]    (J)    at (-4.2,  0.0) {$J(x)$};
\node[box, fill=cPrime, draw=dPrime]    (pi)   at ( 0.0,  0.0) {$\pi(x)$};
\node[box, fill=cPrime, draw=dPrime]    (th)   at ( 5.0,  0.0) {$\theta(x)$};
\node[box, fill=cPrime, draw=dPrime]    (ps)   at ( 9.7,  0.0) {$\psi(x)$};

% 第二层：积性函数 (y=-3.0)
% 节点定义修改
\node[box, fill=cMu, draw=dMu] (mu)  at (0.3, -3.0) {$\mu(n)$};
\node[box, fill=cMu, draw=dMu] (Lam) at (5.0, -3.0) {$\Lambda(n)$};

%Weierstrass乘积
\node[box, fill=cProd, draw=dProd, text=tProd]  (prod)   at ( 9.5,  -17.0) {\tiny\text{Hadamard 乘积}};
% M(x) 节点已移除：初等等价信息直接标注在 RH 框内

% 第三层：Li 与 Gamma (y=-6.0 至 -7.0)
\node[box, fill=cLi, draw=dLi] (Li) at (-2.0, -6.0) {$\Li(x)$};

% Gamma(s) 上下分色：上部标收敛域，下部标延拓
\node[splitbox, rectangle split draw splits=false, minimum width=5cm,
      minimum height=4.7cm,
      rectangle split part fill={cGammaU, cGammaL}] (gamma) at (7.0, -6.0)
    {
     \centering \color{tGamma}$\operatorname{Re}(s)>0$
     \nodepart{two}
     \raisebox{6pt}{\color{tGamma}$\mathbb{C}\setminus\{0,-1,-2,\dots\}$}
    };
% Gamma(s) 标注（固定在节点中线左侧）
\node[font=\small, left=-30pt, text=tGamma] at (gamma.west) {$\Gamma(s)$};

% 【修订3】ζ(s) 节点：下半部分补全函数方程左侧 ζ(s)=
\node[splitbox, rectangle split draw splits=false, minimum height=4cm,
      rectangle split part fill={cZetaU, cZetaL}] (zeta) at (0.3, -10.5)
    {\centering\small
     \color{tZeta}$\displaystyle
       \sum_{n=1}^{\infty}\frac{1}{n^s}
       =\prod_{p}\dfrac{1}{1-p^{-s}}
       \quad(\operatorname{Re}(s)>1)$%
     \vphantom{$\displaystyle\sum_{n=1}^{\infty}$}%
     \nodepart{two}
     % 修订：补全 ζ(s)= 使函数方程完整
     \textcolor{tZeta}{$\displaystyle\zeta(s)=
       2^s\pi^{s-1}
       \sin\!\left(\frac{\pi s}{2}\right)\Gamma(1-s)\,\zeta(1-s)$%
     \vphantom{$\displaystyle\sum_{n=1}^{\infty}\mathbb{C}\setminus\{1\}$}}
    };
% ζ(s) 标注


% ξ(s) 节点：下移至与 RH 同一水平，右侧对称于 M(x)
\node[box, fill=cXi, draw=dXi, text=tXi, minimum height=1.2cm]
  (xi) at (4.85, -17.0)
  {$\xi(s)$\\ \tiny 非平凡零点$\rho$与$\zeta(s)$相同\\
  \tiny $\xi(s)=\xi(1-s)$(函数方程推论)
  };


% 【修订8】N(T) 节点：置于 ζ 与 RH 正中间，构成竖直逻辑链
\node[box, fill=cNT, draw=dNT, text=tNT, minimum height=1.3cm]
  (NT) at (2.3, -14.0)
  {$N(T)$\\ \tiny $\zeta(s)$ 零点计数函数\\[1pt]\tiny（幅角原理）};

% ==================== 黎曼猜想节点 ====================
% RH 下移至 y=-20.5，为 N(T) 留出充分空间
\node[rhbox, minimum width=4.6cm, minimum height=2.2cm] (rh) at (-1.9, -17.0) {
    {\tiny\color{tRH} 黎曼猜想 (Riemann Hypothesis, 1859)}\\[3pt]
    {\tiny\color{tRH}$\zeta(s)$ 的所有非平凡零点 $\rho$ 满足
    $\operatorname{Re}(\rho)=\dfrac{1}{2}$}\\[4pt]
{\tiny\color{tRH} 初等等价（梅滕斯函数）：%
  $\displaystyle M(x)=\sum_{n\leq x}\mu(n)$}\\
{\tiny\color{tRH}
  $\mathrm{RH} \iff M(x) = O(x^{1/2+\varepsilon}),\ \forall \varepsilon>0$}
};


% ==================== 连线 ====================

% ----- 第一层内部 -----
\draw[barr, eArith] (J.east) -- (pi.west)
  node[lab, midway, above=14pt, xshift=-2pt, eArith]
    {$\displaystyle\pi(x)=\sum_{n=1}^{\infty}\frac{\mu(n)}{n}J(x^{1/n})$}
  node[lab, midway, below=14pt, xshift=-2pt, eArith]
    {$\displaystyle J(x)=\sum_{n=1}^{\infty}\frac{1}{n}\pi(x^{1/n})$};

\draw[barr, eElem] (pi.east) -- (th.west)
  node[lab, midway, above=14pt, eElem]
    {$\displaystyle\theta(x)=\pi(x)\log x-\int_2^x\frac{\pi(t)}{t}\,\mathrm{d}t$}
  node[lab, midway, above={14pt}, yshift={-50pt}, eElem]
    {\tiny $\displaystyle\pi(x) = \dfrac{\theta(x)}{\log x}+ \int_{2}^{x} \frac{\theta(t)}{t\,(\log t)^2} \,\mathrm{d}t$};

\draw[barr, eElem] (th.east) -- (ps.west)
  node[lab, midway, above=11pt, xshift={5pt}, eElem]
    {$\displaystyle\psi(x)=\sum_{k=1}^{\infty}\theta(x^{1/k})$}
  node[lab, midway, above=-31pt, xshift={5pt}, yshift={-8pt}, eElem]
    {\tiny $\displaystyle\theta(x)=\sum_{k=1}^{\infty}\mu(k)\psi(x^{1/k})$};


% ----- 第二层内部 -----
\draw[barr, eArith] (mu.east) -- (Lam.west)
  node[lab, midway, yshift=20, xshift=0.3cm, eArith]
    {$\displaystyle\Lambda(n)=-\sum_{d\mid n}\mu(d)\log d$
     \quad{\rm (M\"{o}bius 反演)}}
  node[lab, midway, yshift={-12pt-0.5cm}, xshift=0, eArith]
    {$\displaystyle\log n=\sum_{d\mid n}\Lambda(d)$};

% μ(n) → M(x) 的定义连线已省去（见节点内文字），避免与已有连线交叉

% Lambda -> psi
\draw[arr, eDef] (Lam.east) to[out=0, in=-90, looseness=1.0]
  ($(ps.south west)+(30pt,0)$)
  node[lab, pos=0.5, yshift=-90, xshift=250, eDef]
    {$\displaystyle\psi(x)=\sum_{n\le x}\Lambda(n)$};

% ----- 第三层：Li 和 Gamma -----

% 【修订5】Li(x) → π(x)：改为渐近等价语义（非因果推导）
\draw[arr, ePNT]
  ($(Li.north)!.75!(Li.north east)$)
  to[out=90, in=-135, looseness=1.8]
  ($(pi.south west)+(35pt,0)$)
  node[lab, pos=0.4, yshift=-45, xshift=45, ePNT, align=left]
    {\textbf{素数定理（PNT）}：$\pi(x)\sim\Li(x)$\\
     (若RH成立，$\pi(x)=\Li(x)+O(\sqrt{x}\log x)$)};

% Li -> J（显式公式主项；属 eExpl：第2/9章灰白，第15章全文总图着色）
\draw[arr, eExpl]
  (Li.north) to[out=135, in=-90, looseness=1.5]
  ($(J.south east)+(-22pt,0)$)
  node[lab, pos=0.5, yshift=-115, xshift=-65, eExpl, align=center]
    {$\Li(x)$ 是 $J(x)$ 显式\\[2pt]公式的主项};

% Gamma（Re>0）→ ζ：Euler 积分表示（该积分公式本身要求 Re(s)>1）
\draw[arr, eAnal]
  ([yshift=6pt] gamma.west)
  to[out=180, in=90, looseness=0.8]
  ($(zeta.north)+(30pt,0)$)
  node[lab, pos=0.5, yshift=-235, xshift=90,eAnal, align=center]
    {$\displaystyle\zeta(s)\Gamma(s)
      =\int_0^{\infty}\frac{x^{s-1}}{e^x-1}\,\mathrm{d}x$\\
     $(\operatorname{Re}(s)>1)$};

% μ(n) → ζ(s)
\draw[arr, eAnal]
  (mu.south) to[out=-90, in=90, looseness=1.2]
  ($(zeta.north)+({-20pt+0.3cm},0)$)
  node[lab, pos=0.5, yshift={-120pt-0.5cm-0.4cm}, xshift={45pt-5pt}, eAnal]
    {$\displaystyle\frac{1}{\zeta(s)}=\sum_{n=1}^{\infty}\frac{\mu(n)}{n^s}$\\
     $(\operatorname{Re}(s)>1)$};

% ----- J(x) 与 ψ(x) 的 Stieltjes 积分互逆关系 -----
% 【修订4】加注 dψ(t)/ln t 在 t=1 附近需 t>1 的说明
\draw[barr, eArith, line width=0.8pt, rounded corners=8pt]
  ($(J.north)+(0,0.0cm)$)
  -- ++(0,2.0cm)
  -- ++(13.5cm,0)
  -- ++(0,-2.0cm)
  node[lab, pos=0.35, above=6pt, yshift={10pt+0.5cm}, xshift=-200, eArith, align=center]
    {$\displaystyle\psi(x)=\int_{2^-}^x\log t\,\mathrm{d}J(t)$\quad(Stieltjes 积分，$J\to\psi$)}
  node[lab, pos=0.65, below=6pt, yshift={32pt+0.2cm}, xshift=-205, eArith, align=center]
    {$\displaystyle J(x)=\int_{2^-}^x\frac{\mathrm{d}\psi(t)}{\log t}$
     \quad(Stieltjes 反演，$\psi\to J$；积分自 $t>1$)};

% ----- 第四层：ζ 和 ξ 与第二层的连线 -----

% ζ → ψ：显式公式（Perron 积分 + 留数，单向；x≠p^k 时公式精确成立）
\draw[arr, eExpl, rounded corners=30pt]
  ($(zeta.east)+(0, -10pt)$) -| ($(ps.south)+(30pt,0)$)
  node[lab, pos=0.85, right=10pt, yshift=-170, xshift={-170pt-0.2cm},eExpl, align=left]
{$\psi(x)$ 显式公式（Perron 积分，$\zeta\to\psi$）\\
     $\displaystyle \psi(x) = \frac{1}{2\pi i}\int_{c-i\infty}^{c+i\infty}\!\left(-\frac{\zeta'(s)}{\zeta(s)}\right)\!\frac{x^s}{s}\,\mathrm{d}s$\\
     $\displaystyle = x-\sum_{\rho}\frac{x^{\rho}}{\rho}-\log(2\pi)-\tfrac{1}{2}\log(1-x^{-2})$\\
     \tiny $x>1,\;x\neq p^k$；$\rho$ 按 $|\mathrm{Im}(\rho)|$ 升序配对，条件收敛};

% ψ → ζ：Mellin 变换（单向）
\draw[arr, eExpl, rounded corners=30pt]
  ($(ps.south)+(40pt,0)$) |- ($(zeta.east)+(0, -20pt)$)
  node[lab, pos=0.8, yshift=-27, xshift=30,eExpl, align=center]
    {Mellin 变换（$\psi(x) \longleftrightarrow -\dfrac{\zeta'(s)}{\zeta(s)}$，即负对数导数）\\
     $\displaystyle -\frac{\zeta'(s)}{\zeta(s)} = s\int_{1}^{\infty}\frac{\psi(x)}{x^{s+1}}\,\mathrm{d}x$};
       
       
       
       
       
       
       
% Lambda -> zeta（对数微分 Dirichlet 级数）
\draw[arr, eAnal]
  ($(Lam.south west)+(6pt,0)$) to[out=-90, in=90, looseness=0.4]
  ($(zeta.north west)!0.50!(zeta.north east)$)
  node[lab, pos=0.5, yshift={-120pt-0.5cm-0.4cm+1cm}, xshift={145pt-5pt}, eAnal]
    {$\displaystyle-\frac{\zeta'(s)}{\zeta(s)}
      =\sum_{n=1}^{\infty}\frac{\Lambda(n)}{n^s}$};

% J(x) 显式公式（ζ → J）
\draw[arr, eExpl, rounded corners=50pt]
  ($(zeta.south west)+(0,6pt)$) -- ++(-60pt,0) -- ($(J.south west)+(5pt,0)$)
  node[lab, pos=0.5, yshift=-203, xshift={75pt-0.2cm}, eExpl, align=center]
{\textbf{Perron 公式} $\Rightarrow$ $J(x)$ 显式展开\\
     $\displaystyle
     \begin{aligned}
         &J(x)=\frac{1}{2\pi i}\int_{c-i\infty}^{c+i\infty}
           \frac{\log\zeta(s)}{s}\,x^s\,\mathrm{d}s\\
         &=\Li(x)-\sum_{\rho}\Li(x^{\rho})
           -\log 2+\int_x^{\infty}\frac{\mathrm{d}t}{t(t^2-1)\log t}
     \end{aligned}$\\
     \tiny $x>1,\;x\neq p^k$；$\rho$ 按 $|\mathrm{Im}(\rho)|$ 升序配对，条件收敛};

% ζ → J：Mellin 变换（ln ζ）
\draw[arr, eAnal, rounded corners=5pt]
  (zeta.north west) to[out=110, in=-90, looseness=1.5] (J.south)
  node[lab, pos=0.4, yshift=-205, xshift=-55, eAnal, align=center]
    {$\displaystyle\log\zeta(s)=s\int_1^{\infty}J(x)\,x^{-s-1}\,\mathrm{d}x$\\
     $(\operatorname{Re}(s)>1)$};

% ζ → ξ（定义）：ξ 已下移，连线沿右侧空白区域向右下延伸
\draw[arr, eLate, rounded corners=15pt]
  ($(zeta.south east)!0.15!(zeta.north east)$)
  -- ++(1.0, 0)
  -- ++(0, -3.8)
  -- ($(xi.north)+(0pt,0)$)
  node[lab, pos=0.25, right=4pt,yshift=30, xshift=0, eLate, align=center]
    {$\xi(s)=\dfrac{1}{2}s(s-1)\pi^{-s/2}\Gamma(s/2)\zeta(s)$  
};






% xi -> prod (Weierstrass product connection)
\draw[arr, eLate]
  %(xi.south) -- (prod.south)
  (xi.east) --(prod.west)
  node[midway, right=4pt, yshift=28, xshift=-45, eLate]
    {\tiny$\displaystyle\xi(s)=e^{A+Bs}\prod_{\rho}\left(1-\frac{s}{\rho}\right)e^{s/\rho}$};

% prod -> ps (zero sum in explicit formula)
%\draw[arr, violet!70!black, rounded corners=15pt]
  %(prod.east) -- ++(4.4, 0) |- (ps.south)
  %node[lab, pos=0.6, right=2pt, violet!70!black, align=left]
   % {对数微分与 Perron 积分\\
    % 给出显式公式中\\
    % 零点和 $\sum_{\rho}\frac{x_{\rho}}{\rho}$ 项};
%
% ζ 自身左侧弧线（解析延拓）：与表 tab:spectrum-stage8-guide 一致
\draw[arr, eExt, rounded corners=5pt]
  ($(zeta.north west)!0.15!(zeta.south west)$)
  to[out=180, in=180, looseness=1.6]
  ($(zeta.north west)!0.85!(zeta.south west)$);
% 标签用独立节点锚定在 ζ 西侧弧线旁（避免 path 末 node 漂移）
% 单行；中心在 zeta.west 右 2.5cm（相对初版再累计右移）
\node[lab, eExt, align=center, anchor=center]
  at ([xshift=2.5cm]zeta.west)
  {$\zeta$ 亚纯延拓至 $\mathbb{C}\setminus\{1\}$（Hankel / $\vartheta$ 两法）};

% ζ → N(T)：竖直向下（弧线）
\draw[arr, eZero]
  ($(zeta.south east)!0.10!(zeta.north east)$)
  to[out=0, in=90] ++(0.8, -0.8)
  to[out=-90, in=0] (NT.east)
  node[lab, pos=0.6, right=6pt, yshift=-357, xshift=5, eZero, align=center]
    {$\displaystyle N(T)=\frac{T}{2\pi}\log\frac{T}{2\pi e}+O(\log T)$\\
     \tiny 零点计数公式（von Mangoldt）};

% 【已删去】ζ → RH 蓝色直连线（设计意图改为ζ→N(T)→RH线性链）
%\draw[arr, eArith, rounded corners=10pt]
%  (zeta.south)
%  -- ++(0, -3.35)
%  -- ++(-4.1, 0)
%  -- ($(rh.north)!0.80!(rh.north west)$)
%  node[lab, pos=0.35, left=6pt, yshift=20pt,xshift=120,eArith, align=center]
%    {
%     $\zeta(s)$ 的非平凡零点全部位于直线 $\operatorname{Re}(s)=\dfrac12$ 上};
% RH -- P(s) 画线 + 文字
%\draw[arr, cyan!60!black]
 % ($(rh.east)!0.60!(rh.south east)$) -- (prod.west)    % 起点 -- 终点（必须有！）
  %node[midway, right=6pt, font=\tiny, cyan!60!black, align=center]
  %  {$\zeta(s) = \frac{e^{s(\log(2\pi) - 1)}}{2(s-1)} \prod_{n=1}^{\infty} \left(1 + \frac{s}{2n}\right)e^{-s/2n} P(s)$};

% ξ(s) ↔ RH、N(T) → RH：虚线标注（与 RH 相关的谱系连线）

% N(T) → RH：完成 ζ→N(T)→RH 线性逻辑链（虚线）
\draw[arr, eZero, cond, rounded corners=8pt]
  (NT.south west)
  to[out=-150, in=60, looseness=0.9]
  ($(rh.north)+(20pt, 0)$)
node[lab, pos=0.55, left=4pt, yshift=-430, xshift=-25, eZero, align=center]
    {$\mathrm{RH} \iff N_0(T)=N(T),\ \forall\,T>0$\\
   \tiny（$N_0(T)$：临界线 $\operatorname{Re}(s)=\dfrac{1}{2}$ 上虚部 $\leq T$ 的零点数）};

\draw[barr, eZero, cond]
  (xi.west) -- ($(rh.east)!0.0!(rh.south east)$)
  node[lab, midway, above=20pt,yshift=3pt,xshift=5, eZero, align=center]
    {$\mathrm{RH}\iff\xi\!\left(\dfrac{1}{2}+it\right)$\\
     的所有零点 $t$ 均为实数};

% M(x)→RH 箭头已删去：等价信息直接写在 RH 框内

% Gamma 自身解析延拓弧线：与表 tab:spectrum-stage8-guide 一致
\draw[arr, eExt, rounded corners=1.5pt]
  ($(gamma.north east)+(0,-5pt)$)
  to[out=0, in=0, looseness=1.6]
  ($(gamma.south east)+(0,5pt)$);
% 标签用独立节点锚定在 Γ 东侧弧线旁（仅递推式；延拓域已在框下半）
% 相对初版 (gamma.east+14pt)：净向左 3.7cm，上移 2pt
\node[lab, eExt, align=center, anchor=west]
  at ($(gamma.east)+({14pt-3.7cm},2pt)$)
  {递推延拓 $\Gamma(s)=\Gamma(s+1)/s$};

% Gamma（延拓后）→ ζ：Hankel 围道
\draw[arr, eExt]
  ($(gamma.south west)+(40pt,0)$)
  to[out=180, in=0, looseness=1.0]
  ($(zeta.south east)!0.45!(zeta.north east)$)
  node[lab, pos=0.5, yshift=-210, xshift=215, eExt, align=center]
    {$\displaystyle\zeta(s) = -\frac{\Gamma(1-s)}{2\pi i}\oint_{\mathrm{H}}\frac{(-z)^{s-1}}{e^z-1}\,\mathrm{d}z$\\
     解析延拓到 $\mathbb{C}\setminus\{1\}$，Hankel 围道消去分支切割};

\end{scope}% 主图

% ==================== 图例（取自 A-2 详细文案 + 分阶段灰白色） ====================
\node[draw=black!40, fill=gray!8, rounded corners=8pt, thick, align=center, inner sep=10pt, font=\tiny]
  (legend) at (0, -23.5) {
  \begin{tabular}{@{}ll@{\hspace{1.0cm}}ll@{}}
    \multicolumn{2}{l}{\scriptsize\textbf{节点颜色分类图例}} &
    \multicolumn{2}{l}{\scriptsize\textbf{知识谱系连线语义图例}} \\[6pt]
    \tikz[baseline=-0.5ex]\fill[lgPrime, draw=black!70, thick, rounded corners=1.5pt] (0,-0.12) rectangle (0.55,0.12); &
    \textcolor{lgTprime}{\textbf{素数计数函数}（实变，如 $\pi(x)$）} &
    \textcolor{lgEArith}{\rule[0.5ex]{16pt}{1.4pt}} &
    \textcolor{lgTarith}{\textbf{算术互逆}（M\"{o}bius / Stieltjes 反演）} \\[3pt]
    \tikz[baseline=-0.5ex]\fill[lgMu, draw=black!70, thick, rounded corners=1.5pt] (0,-0.12) rectangle (0.55,0.12); &
    \textcolor{lgTmu}{\textbf{算术积性函数}（离散，如 $\mu(n)$）} &
    \textcolor{lgEDef}{\rule[0.5ex]{16pt}{1.4pt}} &
    \textcolor{lgTdef}{\textbf{对象定义}（构造定义与主项提取）} \\[3pt]
    \tikz[baseline=-0.5ex]\fill[lgLi, draw=black!70, thick, rounded corners=1.5pt] (0,-0.12) rectangle (0.55,0.12); &
    \textcolor{lgTli}{\textbf{渐近逼近主项}（如 $\Li(x)$）} &
    \textcolor{lgEElem}{\rule[0.5ex]{16pt}{1.4pt}} &
    \textcolor{lgTelem}{\textbf{初等计数变换}（数论函数基本恒等式）} \\[3pt]
    \tikz[baseline=-0.5ex]\fill[lgXi, draw=black!55, thick, rounded corners=1.5pt] (0,-0.12) rectangle (0.55,0.12); &
    \textcolor{lgTxi}{\textbf{辅助复变函数}（如 $\xi(s)$）} &
    \textcolor{lgEZero}{\rule[0.5ex]{16pt}{1.4pt}} &
    \textcolor{lgTzero}{\textbf{零点结构与猜想}（分布特征及等价命题）} \\[3pt]
    \tikz[baseline=-0.5ex]\fill[lgRH, draw=black!55, thick, rounded corners=1.5pt] (0,-0.12) rectangle (0.55,0.12); &
    \textcolor{lgTrh}{\textbf{核心猜想与目标}（如 RH）} &
    \textcolor{lgEAnal}{\rule[0.5ex]{16pt}{1.4pt}} &
    \textcolor{lgTanal}{\textbf{分析桥梁}（积分变换与 Dirichlet 级数）} \\[3pt]
    \tikz[baseline=-0.5ex]{\clip[rounded corners=1.5pt] (0,-0.12) rectangle (0.55,0.12); \fill[lgGL] (0,-0.12) rectangle (0.55,0); \fill[lgGU] (0,0) rectangle (0.55,0.12); \draw[black!55, thick, rounded corners=1.5pt] (0,-0.12) rectangle (0.55,0.12); \draw[black!55, thick] (0,0) -- (0.55,0);} &
    \textcolor{lgTgamma}{\textbf{解析延拓复变}（初始域/延拓域）} &
    \textcolor{lgEPNT}{\rule[0.5ex]{16pt}{1.4pt}} &
    \textcolor{lgTpnt}{\textbf{渐近定理}（PNT 等确立的渐近推论）} \\[3pt]
    & & \textcolor{lgEExt}{\rule[0.5ex]{16pt}{1.4pt}} &
    \textcolor{lgText}{\textbf{解析延拓}（复围道积分与全纯延拓）} \\[3pt]
    & & \tikz[baseline=-0.5ex]{\draw[lgEZero, line width=1.2pt] (0.00,0) -- (0.12,0); \draw[lgEZero, line width=1.2pt] (0.17,0) -- (0.29,0); \draw[lgEZero, line width=1.2pt] (0.34,0) -- (0.46,0);} &
    \textbf{未证}
  \end{tabular}
};

\end{scope}% 位移：左 1.0cm、上 2.0cm

\end{tikzpicture}

% 几何与 黎曼猜想知识谱系关系图-A-2.tex 同步；未涉及者仅灰白化（含图例）
%
% \spectrumStage 仅为内部进度码（历史命名），不是章号：
%   2  → 第一次着色：第2章 初等算术层（fig:spectrum-stage2）
%   8  → 第二次着色：第9章 延拓/函数方程（fig:spectrum-stage8）
%  16  → 第三次着色：第15章 全文总图（fig:function-relations）
% 正文表述用「第一次 / 第二次 / 第三次着色」或章标签，勿写「第8章图」「第16章图」。
\providecommand{\spectrumStage}{16}
\makeatletter
\colorlet{cDim}{gray!16}\colorlet{dDim}{gray!48}\colorlet{tDim}{gray!52}
\ifnum\spectrumStage=2
  \colorlet{cPrime}{blue!12}\colorlet{dPrime}{black}\colorlet{tPrime}{black}
  \colorlet{cMu}{teal!15}\colorlet{dMu}{teal!60!black}\colorlet{tMu}{black}
  \colorlet{cLi}{yellow!25}\colorlet{dLi}{orange!70!black}\colorlet{tLi}{black}
  \colorlet{cGammaU}{gray!16}\colorlet{cGammaL}{gray!16}\colorlet{dGamma}{gray!48}\colorlet{tGamma}{gray!52}
  \colorlet{cZetaU}{gray!16}\colorlet{cZetaL}{gray!16}\colorlet{tZeta}{gray!52}
  \colorlet{cXi}{gray!16}\colorlet{dXi}{gray!48}\colorlet{tXi}{gray!52}
  \colorlet{cNT}{gray!16}\colorlet{dNT}{gray!48}\colorlet{tNT}{gray!52}
  \colorlet{cProd}{gray!16}\colorlet{dProd}{gray!48}\colorlet{tProd}{gray!52}
  \colorlet{cRH}{gray!16}\colorlet{dRH}{gray!48}\colorlet{tRH}{gray!52}
  \colorlet{eArith}{blue!70!black}\colorlet{eElem}{violet!80!black}\colorlet{eDef}{black!80}
  \colorlet{ePNT}{brown!70!black}\colorlet{eAnal}{gray!42}\colorlet{eExpl}{gray!42}
  \colorlet{eExt}{gray!42}\colorlet{eZero}{gray!42}\colorlet{eLate}{gray!42}
  \colorlet{lgPrime}{blue!12}\colorlet{lgMu}{teal!15}\colorlet{lgLi}{yellow!25}
  \colorlet{lgXi}{gray!22}\colorlet{lgRH}{gray!22}\colorlet{lgGU}{gray!22}\colorlet{lgGL}{gray!22}
  \colorlet{lgEArith}{blue!70!black}\colorlet{lgEDef}{black!80}\colorlet{lgEElem}{violet!80!black}
  \colorlet{lgEZero}{gray!45}\colorlet{lgEAnal}{gray!45}\colorlet{lgEPNT}{brown!70!black}\colorlet{lgEExt}{gray!45}
  \colorlet{lgTprime}{black}\colorlet{lgTmu}{black}\colorlet{lgTli}{black}
  \colorlet{lgTxi}{gray!50}\colorlet{lgTrh}{gray!50}\colorlet{lgTgamma}{gray!50}
  \colorlet{lgTarith}{black}\colorlet{lgTdef}{black}\colorlet{lgTelem}{black}
  \colorlet{lgTzero}{gray!50}\colorlet{lgTanal}{gray!50}\colorlet{lgTpnt}{black}\colorlet{lgText}{gray!50}
\else
\ifnum\spectrumStage=8
  \colorlet{cPrime}{blue!12}\colorlet{dPrime}{black}\colorlet{tPrime}{black}
  \colorlet{cMu}{teal!15}\colorlet{dMu}{teal!60!black}\colorlet{tMu}{black}
  \colorlet{cLi}{yellow!25}\colorlet{dLi}{orange!70!black}\colorlet{tLi}{black}
  \colorlet{cGammaU}{blue!15}\colorlet{cGammaL}{red!15}\colorlet{dGamma}{black}\colorlet{tGamma}{black}
  \colorlet{cZetaU}{blue!15}\colorlet{cZetaL}{red!15}\colorlet{tZeta}{black}
  \colorlet{cXi}{gray!16}\colorlet{dXi}{gray!48}\colorlet{tXi}{gray!52}
  \colorlet{cNT}{gray!16}\colorlet{dNT}{gray!48}\colorlet{tNT}{gray!52}
  \colorlet{cProd}{gray!16}\colorlet{dProd}{gray!48}\colorlet{tProd}{gray!52}
  \colorlet{cRH}{gray!16}\colorlet{dRH}{gray!48}\colorlet{tRH}{gray!52}
  \colorlet{eArith}{blue!70!black}\colorlet{eElem}{violet!80!black}\colorlet{eDef}{black!80}
  \colorlet{ePNT}{brown!70!black}\colorlet{eAnal}{green!50!black}\colorlet{eExpl}{gray!42}
  \colorlet{eExt}{red!70!black}\colorlet{eZero}{gray!42}\colorlet{eLate}{gray!42}
  \colorlet{lgPrime}{blue!12}\colorlet{lgMu}{teal!15}\colorlet{lgLi}{yellow!25}
  \colorlet{lgXi}{gray!22}\colorlet{lgRH}{gray!22}\colorlet{lgGU}{blue!15}\colorlet{lgGL}{red!15}
  \colorlet{lgEArith}{blue!70!black}\colorlet{lgEDef}{black!80}\colorlet{lgEElem}{violet!80!black}
  \colorlet{lgEZero}{gray!45}\colorlet{lgEAnal}{green!50!black}\colorlet{lgEPNT}{brown!70!black}\colorlet{lgEExt}{red!70!black}
  \colorlet{lgTprime}{black}\colorlet{lgTmu}{black}\colorlet{lgTli}{black}
  \colorlet{lgTxi}{gray!50}\colorlet{lgTrh}{gray!50}\colorlet{lgTgamma}{black}
  \colorlet{lgTarith}{black}\colorlet{lgTdef}{black}\colorlet{lgTelem}{black}
  \colorlet{lgTzero}{gray!50}\colorlet{lgTanal}{black}\colorlet{lgTpnt}{black}\colorlet{lgText}{black}
\else
  \colorlet{cPrime}{blue!12}\colorlet{dPrime}{black}\colorlet{tPrime}{black}
  \colorlet{cMu}{teal!15}\colorlet{dMu}{teal!60!black}\colorlet{tMu}{black}
  \colorlet{cLi}{yellow!25}\colorlet{dLi}{orange!70!black}\colorlet{tLi}{black}
  \colorlet{cGammaU}{blue!15}\colorlet{cGammaL}{red!15}\colorlet{dGamma}{black}\colorlet{tGamma}{black}
  \colorlet{cZetaU}{blue!15}\colorlet{cZetaL}{red!15}\colorlet{tZeta}{black}
  \colorlet{cXi}{cyan!15}\colorlet{dXi}{cyan!60!black}\colorlet{tXi}{black}
  \colorlet{cNT}{cyan!15}\colorlet{dNT}{cyan!60!black}\colorlet{tNT}{black}
  \colorlet{cProd}{cyan!15}\colorlet{dProd}{cyan!60!black}\colorlet{tProd}{black}
  \colorlet{cRH}{orange!15}\colorlet{dRH}{orange!70!black}\colorlet{tRH}{black}
  \colorlet{eArith}{blue!70!black}\colorlet{eElem}{violet!80!black}\colorlet{eDef}{black!80}
  \colorlet{ePNT}{brown!70!black}\colorlet{eAnal}{green!50!black}\colorlet{eExpl}{green!50!black}
  \colorlet{eExt}{red!70!black}\colorlet{eZero}{orange!80!black}\colorlet{eLate}{black!80}
  \colorlet{lgPrime}{blue!12}\colorlet{lgMu}{teal!15}\colorlet{lgLi}{yellow!25}
  \colorlet{lgXi}{cyan!15}\colorlet{lgRH}{orange!15}\colorlet{lgGU}{blue!15}\colorlet{lgGL}{red!15}
  \colorlet{lgEArith}{blue!70!black}\colorlet{lgEDef}{black!80}\colorlet{lgEElem}{violet!80!black}
  \colorlet{lgEZero}{orange!80!black}\colorlet{lgEAnal}{green!50!black}\colorlet{lgEPNT}{brown!70!black}\colorlet{lgEExt}{red!70!black}
  \colorlet{lgTprime}{black}\colorlet{lgTmu}{black}\colorlet{lgTli}{black}
  \colorlet{lgTxi}{black}\colorlet{lgTrh}{black}\colorlet{lgTgamma}{black}
  \colorlet{lgTarith}{black}\colorlet{lgTdef}{black}\colorlet{lgTelem}{black}
  \colorlet{lgTzero}{black}\colorlet{lgTanal}{black}\colorlet{lgTpnt}{black}\colorlet{lgText}{black}
\fi\fi
\makeatother

\begin{tikzpicture}[scale=1.0, transform shape,
  box/.style    = {draw, rounded corners=5pt, minimum width=2.8cm,
                   minimum height=0.9cm, align=center, font=\small, thick},
  zbox/.style   = {draw, rounded corners=5pt, minimum width=2.8cm,
                   minimum height=0.9cm, align=center, font=\small,
                   thick, fill=gray!18},
  splitbox/.style = {draw, rounded corners=5pt, minimum width=7cm,
                     minimum height=2.2cm, align=center, font=\small, thick,
                     rectangle split, rectangle split parts=2,
                     rectangle split part fill={red!8, white}},
  arr/.style    = {-{Stealth[length=6pt,width=5pt]}, thick},
  barr/.style   = {{Stealth[length=6pt,width=5pt]}-{Stealth[length=6pt,width=5pt]}, thick},
  % 线型：虚线仅用于 N(T)→RH、ξ↔RH；其余为实线
  proven/.style = {solid},
  cond/.style   = {dashed},
  lab/.style    = {font=\tiny, align=center, inner sep=3pt},
  rhbox/.style  = {draw=dRH, rounded corners=5pt, minimum width=8cm,
                   minimum height=1.0cm, align=center, font=\small,
                   fill=cRH, text=tRH, thick}
]

% 布局与 A-2 同步：左移 1.0cm、上移 2.0cm
% 左右略放宽，避免 Γ/ζ 自环旁公式标注被裁切
\path[use as bounding box]
  (-10.2, -26.2) rectangle (10.0, 7.4);

\begin{scope}[shift={(-1.0,2.0)}]

\node[font=\bfseries\Large, align=center] at (0, 4.35)
  {黎曼猜想知识谱系关系图};

% 主图单独上移 1cm；标题与图例位置不动
\begin{scope}[shift={(-2.75,-1.2)}]
%% ==================== 节点分层 ====================

% 第一层：素数计数函数 (y=0)
\node[box, fill=cPrime, draw=dPrime]    (J)    at (-4.2,  0.0) {$J(x)$};
\node[box, fill=cPrime, draw=dPrime]    (pi)   at ( 0.0,  0.0) {$\pi(x)$};
\node[box, fill=cPrime, draw=dPrime]    (th)   at ( 5.0,  0.0) {$\theta(x)$};
\node[box, fill=cPrime, draw=dPrime]    (ps)   at ( 9.7,  0.0) {$\psi(x)$};

% 第二层：积性函数 (y=-3.0)
% 节点定义修改
\node[box, fill=cMu, draw=dMu] (mu)  at (0.3, -3.0) {$\mu(n)$};
\node[box, fill=cMu, draw=dMu] (Lam) at (5.0, -3.0) {$\Lambda(n)$};

%Weierstrass乘积
\node[box, fill=cProd, draw=dProd, text=tProd]  (prod)   at ( 9.5,  -17.0) {\tiny\text{Hadamard 乘积}};
% M(x) 节点已移除：初等等价信息直接标注在 RH 框内

% 第三层：Li 与 Gamma (y=-6.0 至 -7.0)
\node[box, fill=cLi, draw=dLi] (Li) at (-2.0, -6.0) {$\Li(x)$};

% Gamma(s) 上下分色：上部标收敛域，下部标延拓
\node[splitbox, rectangle split draw splits=false, minimum width=5cm,
      minimum height=4.7cm,
      rectangle split part fill={cGammaU, cGammaL}] (gamma) at (7.0, -6.0)
    {
     \centering \color{tGamma}$\operatorname{Re}(s)>0$
     \nodepart{two}
     \raisebox{6pt}{\color{tGamma}$\mathbb{C}\setminus\{0,-1,-2,\dots\}$}
    };
% Gamma(s) 标注（固定在节点中线左侧）
\node[font=\small, left=-30pt, text=tGamma] at (gamma.west) {$\Gamma(s)$};

% 【修订3】ζ(s) 节点：下半部分补全函数方程左侧 ζ(s)=
\node[splitbox, rectangle split draw splits=false, minimum height=4cm,
      rectangle split part fill={cZetaU, cZetaL}] (zeta) at (0.3, -10.5)
    {\centering\small
     \color{tZeta}$\displaystyle
       \sum_{n=1}^{\infty}\frac{1}{n^s}
       =\prod_{p}\dfrac{1}{1-p^{-s}}
       \quad(\operatorname{Re}(s)>1)$%
     \vphantom{$\displaystyle\sum_{n=1}^{\infty}$}%
     \nodepart{two}
     % 修订：补全 ζ(s)= 使函数方程完整
     \textcolor{tZeta}{$\displaystyle\zeta(s)=
       2^s\pi^{s-1}
       \sin\!\left(\frac{\pi s}{2}\right)\Gamma(1-s)\,\zeta(1-s)$%
     \vphantom{$\displaystyle\sum_{n=1}^{\infty}\mathbb{C}\setminus\{1\}$}}
    };
% ζ(s) 标注


% ξ(s) 节点：下移至与 RH 同一水平，右侧对称于 M(x)
\node[box, fill=cXi, draw=dXi, text=tXi, minimum height=1.2cm]
  (xi) at (4.85, -17.0)
  {$\xi(s)$\\ \tiny 非平凡零点$\rho$与$\zeta(s)$相同\\
  \tiny $\xi(s)=\xi(1-s)$(函数方程推论)
  };


% 【修订8】N(T) 节点：置于 ζ 与 RH 正中间，构成竖直逻辑链
\node[box, fill=cNT, draw=dNT, text=tNT, minimum height=1.3cm]
  (NT) at (2.3, -14.0)
  {$N(T)$\\ \tiny $\zeta(s)$ 零点计数函数\\[1pt]\tiny（幅角原理）};

% ==================== 黎曼猜想节点 ====================
% RH 下移至 y=-20.5，为 N(T) 留出充分空间
\node[rhbox, minimum width=4.6cm, minimum height=2.2cm] (rh) at (-1.9, -17.0) {
    {\tiny\color{tRH} 黎曼猜想 (Riemann Hypothesis, 1859)}\\[3pt]
    {\tiny\color{tRH}$\zeta(s)$ 的所有非平凡零点 $\rho$ 满足
    $\operatorname{Re}(\rho)=\dfrac{1}{2}$}\\[4pt]
{\tiny\color{tRH} 初等等价（梅滕斯函数）：%
  $\displaystyle M(x)=\sum_{n\leq x}\mu(n)$}\\
{\tiny\color{tRH}
  $\mathrm{RH} \iff M(x) = O(x^{1/2+\varepsilon}),\ \forall \varepsilon>0$}
};


% ==================== 连线 ====================

% ----- 第一层内部 -----
\draw[barr, eArith] (J.east) -- (pi.west)
  node[lab, midway, above=14pt, xshift=-2pt, eArith]
    {$\displaystyle\pi(x)=\sum_{n=1}^{\infty}\frac{\mu(n)}{n}J(x^{1/n})$}
  node[lab, midway, below=14pt, xshift=-2pt, eArith]
    {$\displaystyle J(x)=\sum_{n=1}^{\infty}\frac{1}{n}\pi(x^{1/n})$};

\draw[barr, eElem] (pi.east) -- (th.west)
  node[lab, midway, above=14pt, eElem]
    {$\displaystyle\theta(x)=\pi(x)\log x-\int_2^x\frac{\pi(t)}{t}\,\mathrm{d}t$}
  node[lab, midway, above={14pt}, yshift={-50pt}, eElem]
    {\tiny $\displaystyle\pi(x) = \dfrac{\theta(x)}{\log x}+ \int_{2}^{x} \frac{\theta(t)}{t\,(\log t)^2} \,\mathrm{d}t$};

\draw[barr, eElem] (th.east) -- (ps.west)
  node[lab, midway, above=11pt, xshift={5pt}, eElem]
    {$\displaystyle\psi(x)=\sum_{k=1}^{\infty}\theta(x^{1/k})$}
  node[lab, midway, above=-31pt, xshift={5pt}, yshift={-8pt}, eElem]
    {\tiny $\displaystyle\theta(x)=\sum_{k=1}^{\infty}\mu(k)\psi(x^{1/k})$};


% ----- 第二层内部 -----
\draw[barr, eArith] (mu.east) -- (Lam.west)
  node[lab, midway, yshift=20, xshift=0.3cm, eArith]
    {$\displaystyle\Lambda(n)=-\sum_{d\mid n}\mu(d)\log d$
     \quad{\rm (M\"{o}bius 反演)}}
  node[lab, midway, yshift={-12pt-0.5cm}, xshift=0, eArith]
    {$\displaystyle\log n=\sum_{d\mid n}\Lambda(d)$};

% μ(n) → M(x) 的定义连线已省去（见节点内文字），避免与已有连线交叉

% Lambda -> psi
\draw[arr, eDef] (Lam.east) to[out=0, in=-90, looseness=1.0]
  ($(ps.south west)+(30pt,0)$)
  node[lab, pos=0.5, yshift=-90, xshift=250, eDef]
    {$\displaystyle\psi(x)=\sum_{n\le x}\Lambda(n)$};

% ----- 第三层：Li 和 Gamma -----

% 【修订5】Li(x) → π(x)：改为渐近等价语义（非因果推导）
\draw[arr, ePNT]
  ($(Li.north)!.75!(Li.north east)$)
  to[out=90, in=-135, looseness=1.8]
  ($(pi.south west)+(35pt,0)$)
  node[lab, pos=0.4, yshift=-45, xshift=45, ePNT, align=left]
    {\textbf{素数定理（PNT）}：$\pi(x)\sim\Li(x)$\\
     (若RH成立，$\pi(x)=\Li(x)+O(\sqrt{x}\log x)$)};

% Li -> J（显式公式主项；属 eExpl：第2/9章灰白，第15章全文总图着色）
\draw[arr, eExpl]
  (Li.north) to[out=135, in=-90, looseness=1.5]
  ($(J.south east)+(-22pt,0)$)
  node[lab, pos=0.5, yshift=-115, xshift=-65, eExpl, align=center]
    {$\Li(x)$ 是 $J(x)$ 显式\\[2pt]公式的主项};

% Gamma（Re>0）→ ζ：Euler 积分表示（该积分公式本身要求 Re(s)>1）
\draw[arr, eAnal]
  ([yshift=6pt] gamma.west)
  to[out=180, in=90, looseness=0.8]
  ($(zeta.north)+(30pt,0)$)
  node[lab, pos=0.5, yshift=-235, xshift=90,eAnal, align=center]
    {$\displaystyle\zeta(s)\Gamma(s)
      =\int_0^{\infty}\frac{x^{s-1}}{e^x-1}\,\mathrm{d}x$\\
     $(\operatorname{Re}(s)>1)$};

% μ(n) → ζ(s)
\draw[arr, eAnal]
  (mu.south) to[out=-90, in=90, looseness=1.2]
  ($(zeta.north)+({-20pt+0.3cm},0)$)
  node[lab, pos=0.5, yshift={-120pt-0.5cm-0.4cm}, xshift={45pt-5pt}, eAnal]
    {$\displaystyle\frac{1}{\zeta(s)}=\sum_{n=1}^{\infty}\frac{\mu(n)}{n^s}$\\
     $(\operatorname{Re}(s)>1)$};

% ----- J(x) 与 ψ(x) 的 Stieltjes 积分互逆关系 -----
% 【修订4】加注 dψ(t)/ln t 在 t=1 附近需 t>1 的说明
\draw[barr, eArith, line width=0.8pt, rounded corners=8pt]
  ($(J.north)+(0,0.0cm)$)
  -- ++(0,2.0cm)
  -- ++(13.5cm,0)
  -- ++(0,-2.0cm)
  node[lab, pos=0.35, above=6pt, yshift={10pt+0.5cm}, xshift=-200, eArith, align=center]
    {$\displaystyle\psi(x)=\int_{2^-}^x\log t\,\mathrm{d}J(t)$\quad(Stieltjes 积分，$J\to\psi$)}
  node[lab, pos=0.65, below=6pt, yshift={32pt+0.2cm}, xshift=-205, eArith, align=center]
    {$\displaystyle J(x)=\int_{2^-}^x\frac{\mathrm{d}\psi(t)}{\log t}$
     \quad(Stieltjes 反演，$\psi\to J$；积分自 $t>1$)};

% ----- 第四层：ζ 和 ξ 与第二层的连线 -----

% ζ → ψ：显式公式（Perron 积分 + 留数，单向；x≠p^k 时公式精确成立）
\draw[arr, eExpl, rounded corners=30pt]
  ($(zeta.east)+(0, -10pt)$) -| ($(ps.south)+(30pt,0)$)
  node[lab, pos=0.85, right=10pt, yshift=-170, xshift={-170pt-0.2cm},eExpl, align=left]
{$\psi(x)$ 显式公式（Perron 积分，$\zeta\to\psi$）\\
     $\displaystyle \psi(x) = \frac{1}{2\pi i}\int_{c-i\infty}^{c+i\infty}\!\left(-\frac{\zeta'(s)}{\zeta(s)}\right)\!\frac{x^s}{s}\,\mathrm{d}s$\\
     $\displaystyle = x-\sum_{\rho}\frac{x^{\rho}}{\rho}-\log(2\pi)-\tfrac{1}{2}\log(1-x^{-2})$\\
     \tiny $x>1,\;x\neq p^k$；$\rho$ 按 $|\mathrm{Im}(\rho)|$ 升序配对，条件收敛};

% ψ → ζ：Mellin 变换（单向）
\draw[arr, eExpl, rounded corners=30pt]
  ($(ps.south)+(40pt,0)$) |- ($(zeta.east)+(0, -20pt)$)
  node[lab, pos=0.8, yshift=-27, xshift=30,eExpl, align=center]
    {Mellin 变换（$\psi(x) \longleftrightarrow -\dfrac{\zeta'(s)}{\zeta(s)}$，即负对数导数）\\
     $\displaystyle -\frac{\zeta'(s)}{\zeta(s)} = s\int_{1}^{\infty}\frac{\psi(x)}{x^{s+1}}\,\mathrm{d}x$};
       
       
       
       
       
       
       
% Lambda -> zeta（对数微分 Dirichlet 级数）
\draw[arr, eAnal]
  ($(Lam.south west)+(6pt,0)$) to[out=-90, in=90, looseness=0.4]
  ($(zeta.north west)!0.50!(zeta.north east)$)
  node[lab, pos=0.5, yshift={-120pt-0.5cm-0.4cm+1cm}, xshift={145pt-5pt}, eAnal]
    {$\displaystyle-\frac{\zeta'(s)}{\zeta(s)}
      =\sum_{n=1}^{\infty}\frac{\Lambda(n)}{n^s}$};

% J(x) 显式公式（ζ → J）
\draw[arr, eExpl, rounded corners=50pt]
  ($(zeta.south west)+(0,6pt)$) -- ++(-60pt,0) -- ($(J.south west)+(5pt,0)$)
  node[lab, pos=0.5, yshift=-203, xshift={75pt-0.2cm}, eExpl, align=center]
{\textbf{Perron 公式} $\Rightarrow$ $J(x)$ 显式展开\\
     $\displaystyle
     \begin{aligned}
         &J(x)=\frac{1}{2\pi i}\int_{c-i\infty}^{c+i\infty}
           \frac{\log\zeta(s)}{s}\,x^s\,\mathrm{d}s\\
         &=\Li(x)-\sum_{\rho}\Li(x^{\rho})
           -\log 2+\int_x^{\infty}\frac{\mathrm{d}t}{t(t^2-1)\log t}
     \end{aligned}$\\
     \tiny $x>1,\;x\neq p^k$；$\rho$ 按 $|\mathrm{Im}(\rho)|$ 升序配对，条件收敛};

% ζ → J：Mellin 变换（ln ζ）
\draw[arr, eAnal, rounded corners=5pt]
  (zeta.north west) to[out=110, in=-90, looseness=1.5] (J.south)
  node[lab, pos=0.4, yshift=-205, xshift=-55, eAnal, align=center]
    {$\displaystyle\log\zeta(s)=s\int_1^{\infty}J(x)\,x^{-s-1}\,\mathrm{d}x$\\
     $(\operatorname{Re}(s)>1)$};

% ζ → ξ（定义）：ξ 已下移，连线沿右侧空白区域向右下延伸
\draw[arr, eLate, rounded corners=15pt]
  ($(zeta.south east)!0.15!(zeta.north east)$)
  -- ++(1.0, 0)
  -- ++(0, -3.8)
  -- ($(xi.north)+(0pt,0)$)
  node[lab, pos=0.25, right=4pt,yshift=30, xshift=0, eLate, align=center]
    {$\xi(s)=\dfrac{1}{2}s(s-1)\pi^{-s/2}\Gamma(s/2)\zeta(s)$  
};






% xi -> prod (Weierstrass product connection)
\draw[arr, eLate]
  %(xi.south) -- (prod.south)
  (xi.east) --(prod.west)
  node[midway, right=4pt, yshift=28, xshift=-45, eLate]
    {\tiny$\displaystyle\xi(s)=e^{A+Bs}\prod_{\rho}\left(1-\frac{s}{\rho}\right)e^{s/\rho}$};

% prod -> ps (zero sum in explicit formula)
%\draw[arr, violet!70!black, rounded corners=15pt]
  %(prod.east) -- ++(4.4, 0) |- (ps.south)
  %node[lab, pos=0.6, right=2pt, violet!70!black, align=left]
   % {对数微分与 Perron 积分\\
    % 给出显式公式中\\
    % 零点和 $\sum_{\rho}\frac{x_{\rho}}{\rho}$ 项};
%
% ζ 自身左侧弧线（解析延拓）：与表 tab:spectrum-stage8-guide 一致
\draw[arr, eExt, rounded corners=5pt]
  ($(zeta.north west)!0.15!(zeta.south west)$)
  to[out=180, in=180, looseness=1.6]
  ($(zeta.north west)!0.85!(zeta.south west)$);
% 标签用独立节点锚定在 ζ 西侧弧线旁（避免 path 末 node 漂移）
% 单行；中心在 zeta.west 右 2.5cm（相对初版再累计右移）
\node[lab, eExt, align=center, anchor=center]
  at ([xshift=2.5cm]zeta.west)
  {$\zeta$ 亚纯延拓至 $\mathbb{C}\setminus\{1\}$（Hankel / $\vartheta$ 两法）};

% ζ → N(T)：竖直向下（弧线）
\draw[arr, eZero]
  ($(zeta.south east)!0.10!(zeta.north east)$)
  to[out=0, in=90] ++(0.8, -0.8)
  to[out=-90, in=0] (NT.east)
  node[lab, pos=0.6, right=6pt, yshift=-357, xshift=5, eZero, align=center]
    {$\displaystyle N(T)=\frac{T}{2\pi}\log\frac{T}{2\pi e}+O(\log T)$\\
     \tiny 零点计数公式（von Mangoldt）};

% 【已删去】ζ → RH 蓝色直连线（设计意图改为ζ→N(T)→RH线性链）
%\draw[arr, eArith, rounded corners=10pt]
%  (zeta.south)
%  -- ++(0, -3.35)
%  -- ++(-4.1, 0)
%  -- ($(rh.north)!0.80!(rh.north west)$)
%  node[lab, pos=0.35, left=6pt, yshift=20pt,xshift=120,eArith, align=center]
%    {
%     $\zeta(s)$ 的非平凡零点全部位于直线 $\operatorname{Re}(s)=\dfrac12$ 上};
% RH -- P(s) 画线 + 文字
%\draw[arr, cyan!60!black]
 % ($(rh.east)!0.60!(rh.south east)$) -- (prod.west)    % 起点 -- 终点（必须有！）
  %node[midway, right=6pt, font=\tiny, cyan!60!black, align=center]
  %  {$\zeta(s) = \frac{e^{s(\log(2\pi) - 1)}}{2(s-1)} \prod_{n=1}^{\infty} \left(1 + \frac{s}{2n}\right)e^{-s/2n} P(s)$};

% ξ(s) ↔ RH、N(T) → RH：虚线标注（与 RH 相关的谱系连线）

% N(T) → RH：完成 ζ→N(T)→RH 线性逻辑链（虚线）
\draw[arr, eZero, cond, rounded corners=8pt]
  (NT.south west)
  to[out=-150, in=60, looseness=0.9]
  ($(rh.north)+(20pt, 0)$)
node[lab, pos=0.55, left=4pt, yshift=-430, xshift=-25, eZero, align=center]
    {$\mathrm{RH} \iff N_0(T)=N(T),\ \forall\,T>0$\\
   \tiny（$N_0(T)$：临界线 $\operatorname{Re}(s)=\dfrac{1}{2}$ 上虚部 $\leq T$ 的零点数）};

\draw[barr, eZero, cond]
  (xi.west) -- ($(rh.east)!0.0!(rh.south east)$)
  node[lab, midway, above=20pt,yshift=3pt,xshift=5, eZero, align=center]
    {$\mathrm{RH}\iff\xi\!\left(\dfrac{1}{2}+it\right)$\\
     的所有零点 $t$ 均为实数};

% M(x)→RH 箭头已删去：等价信息直接写在 RH 框内

% Gamma 自身解析延拓弧线：与表 tab:spectrum-stage8-guide 一致
\draw[arr, eExt, rounded corners=1.5pt]
  ($(gamma.north east)+(0,-5pt)$)
  to[out=0, in=0, looseness=1.6]
  ($(gamma.south east)+(0,5pt)$);
% 标签用独立节点锚定在 Γ 东侧弧线旁（仅递推式；延拓域已在框下半）
% 相对初版 (gamma.east+14pt)：净向左 3.7cm，上移 2pt
\node[lab, eExt, align=center, anchor=west]
  at ($(gamma.east)+({14pt-3.7cm},2pt)$)
  {递推延拓 $\Gamma(s)=\Gamma(s+1)/s$};

% Gamma（延拓后）→ ζ：Hankel 围道
\draw[arr, eExt]
  ($(gamma.south west)+(40pt,0)$)
  to[out=180, in=0, looseness=1.0]
  ($(zeta.south east)!0.45!(zeta.north east)$)
  node[lab, pos=0.5, yshift=-210, xshift=215, eExt, align=center]
    {$\displaystyle\zeta(s) = -\frac{\Gamma(1-s)}{2\pi i}\oint_{\mathrm{H}}\frac{(-z)^{s-1}}{e^z-1}\,\mathrm{d}z$\\
     解析延拓到 $\mathbb{C}\setminus\{1\}$，Hankel 围道消去分支切割};

\end{scope}% 主图

% ==================== 图例（取自 A-2 详细文案 + 分阶段灰白色） ====================
\node[draw=black!40, fill=gray!8, rounded corners=8pt, thick, align=center, inner sep=10pt, font=\tiny]
  (legend) at (0, -23.5) {
  \begin{tabular}{@{}ll@{\hspace{1.0cm}}ll@{}}
    \multicolumn{2}{l}{\scriptsize\textbf{节点颜色分类图例}} &
    \multicolumn{2}{l}{\scriptsize\textbf{知识谱系连线语义图例}} \\[6pt]
    \tikz[baseline=-0.5ex]\fill[lgPrime, draw=black!70, thick, rounded corners=1.5pt] (0,-0.12) rectangle (0.55,0.12); &
    \textcolor{lgTprime}{\textbf{素数计数函数}（实变，如 $\pi(x)$）} &
    \textcolor{lgEArith}{\rule[0.5ex]{16pt}{1.4pt}} &
    \textcolor{lgTarith}{\textbf{算术互逆}（M\"{o}bius / Stieltjes 反演）} \\[3pt]
    \tikz[baseline=-0.5ex]\fill[lgMu, draw=black!70, thick, rounded corners=1.5pt] (0,-0.12) rectangle (0.55,0.12); &
    \textcolor{lgTmu}{\textbf{算术积性函数}（离散，如 $\mu(n)$）} &
    \textcolor{lgEDef}{\rule[0.5ex]{16pt}{1.4pt}} &
    \textcolor{lgTdef}{\textbf{对象定义}（构造定义与主项提取）} \\[3pt]
    \tikz[baseline=-0.5ex]\fill[lgLi, draw=black!70, thick, rounded corners=1.5pt] (0,-0.12) rectangle (0.55,0.12); &
    \textcolor{lgTli}{\textbf{渐近逼近主项}（如 $\Li(x)$）} &
    \textcolor{lgEElem}{\rule[0.5ex]{16pt}{1.4pt}} &
    \textcolor{lgTelem}{\textbf{初等计数变换}（数论函数基本恒等式）} \\[3pt]
    \tikz[baseline=-0.5ex]\fill[lgXi, draw=black!55, thick, rounded corners=1.5pt] (0,-0.12) rectangle (0.55,0.12); &
    \textcolor{lgTxi}{\textbf{辅助复变函数}（如 $\xi(s)$）} &
    \textcolor{lgEZero}{\rule[0.5ex]{16pt}{1.4pt}} &
    \textcolor{lgTzero}{\textbf{零点结构与猜想}（分布特征及等价命题）} \\[3pt]
    \tikz[baseline=-0.5ex]\fill[lgRH, draw=black!55, thick, rounded corners=1.5pt] (0,-0.12) rectangle (0.55,0.12); &
    \textcolor{lgTrh}{\textbf{核心猜想与目标}（如 RH）} &
    \textcolor{lgEAnal}{\rule[0.5ex]{16pt}{1.4pt}} &
    \textcolor{lgTanal}{\textbf{分析桥梁}（积分变换与 Dirichlet 级数）} \\[3pt]
    \tikz[baseline=-0.5ex]{\clip[rounded corners=1.5pt] (0,-0.12) rectangle (0.55,0.12); \fill[lgGL] (0,-0.12) rectangle (0.55,0); \fill[lgGU] (0,0) rectangle (0.55,0.12); \draw[black!55, thick, rounded corners=1.5pt] (0,-0.12) rectangle (0.55,0.12); \draw[black!55, thick] (0,0) -- (0.55,0);} &
    \textcolor{lgTgamma}{\textbf{解析延拓复变}（初始域/延拓域）} &
    \textcolor{lgEPNT}{\rule[0.5ex]{16pt}{1.4pt}} &
    \textcolor{lgTpnt}{\textbf{渐近定理}（PNT 等确立的渐近推论）} \\[3pt]
    & & \textcolor{lgEExt}{\rule[0.5ex]{16pt}{1.4pt}} &
    \textcolor{lgText}{\textbf{解析延拓}（复围道积分与全纯延拓）} \\[3pt]
    & & \tikz[baseline=-0.5ex]{\draw[lgEZero, line width=1.2pt] (0.00,0) -- (0.12,0); \draw[lgEZero, line width=1.2pt] (0.17,0) -- (0.29,0); \draw[lgEZero, line width=1.2pt] (0.34,0) -- (0.46,0);} &
    \textbf{未证}
  \end{tabular}
};

\end{scope}% 位移：左 1.0cm、上 2.0cm

\end{tikzpicture}

% 几何与 黎曼猜想知识谱系关系图-A-2.tex 同步；未涉及者仅灰白化（含图例）
%
% \spectrumStage 仅为内部进度码（历史命名），不是章号：
%   2  → 第一次着色：第2章 初等算术层（fig:spectrum-stage2）
%   8  → 第二次着色：第9章 延拓/函数方程（fig:spectrum-stage8）
%  16  → 第三次着色：第15章 全文总图（fig:function-relations）
% 正文表述用「第一次 / 第二次 / 第三次着色」或章标签，勿写「第8章图」「第16章图」。
\providecommand{\spectrumStage}{16}
\makeatletter
\colorlet{cDim}{gray!16}\colorlet{dDim}{gray!48}\colorlet{tDim}{gray!52}
\ifnum\spectrumStage=2
  \colorlet{cPrime}{blue!12}\colorlet{dPrime}{black}\colorlet{tPrime}{black}
  \colorlet{cMu}{teal!15}\colorlet{dMu}{teal!60!black}\colorlet{tMu}{black}
  \colorlet{cLi}{yellow!25}\colorlet{dLi}{orange!70!black}\colorlet{tLi}{black}
  \colorlet{cGammaU}{gray!16}\colorlet{cGammaL}{gray!16}\colorlet{dGamma}{gray!48}\colorlet{tGamma}{gray!52}
  \colorlet{cZetaU}{gray!16}\colorlet{cZetaL}{gray!16}\colorlet{tZeta}{gray!52}
  \colorlet{cXi}{gray!16}\colorlet{dXi}{gray!48}\colorlet{tXi}{gray!52}
  \colorlet{cNT}{gray!16}\colorlet{dNT}{gray!48}\colorlet{tNT}{gray!52}
  \colorlet{cProd}{gray!16}\colorlet{dProd}{gray!48}\colorlet{tProd}{gray!52}
  \colorlet{cRH}{gray!16}\colorlet{dRH}{gray!48}\colorlet{tRH}{gray!52}
  \colorlet{eArith}{blue!70!black}\colorlet{eElem}{violet!80!black}\colorlet{eDef}{black!80}
  \colorlet{ePNT}{brown!70!black}\colorlet{eAnal}{gray!42}\colorlet{eExpl}{gray!42}
  \colorlet{eExt}{gray!42}\colorlet{eZero}{gray!42}\colorlet{eLate}{gray!42}
  \colorlet{lgPrime}{blue!12}\colorlet{lgMu}{teal!15}\colorlet{lgLi}{yellow!25}
  \colorlet{lgXi}{gray!22}\colorlet{lgRH}{gray!22}\colorlet{lgGU}{gray!22}\colorlet{lgGL}{gray!22}
  \colorlet{lgEArith}{blue!70!black}\colorlet{lgEDef}{black!80}\colorlet{lgEElem}{violet!80!black}
  \colorlet{lgEZero}{gray!45}\colorlet{lgEAnal}{gray!45}\colorlet{lgEPNT}{brown!70!black}\colorlet{lgEExt}{gray!45}
  \colorlet{lgTprime}{black}\colorlet{lgTmu}{black}\colorlet{lgTli}{black}
  \colorlet{lgTxi}{gray!50}\colorlet{lgTrh}{gray!50}\colorlet{lgTgamma}{gray!50}
  \colorlet{lgTarith}{black}\colorlet{lgTdef}{black}\colorlet{lgTelem}{black}
  \colorlet{lgTzero}{gray!50}\colorlet{lgTanal}{gray!50}\colorlet{lgTpnt}{black}\colorlet{lgText}{gray!50}
\else
\ifnum\spectrumStage=8
  \colorlet{cPrime}{blue!12}\colorlet{dPrime}{black}\colorlet{tPrime}{black}
  \colorlet{cMu}{teal!15}\colorlet{dMu}{teal!60!black}\colorlet{tMu}{black}
  \colorlet{cLi}{yellow!25}\colorlet{dLi}{orange!70!black}\colorlet{tLi}{black}
  \colorlet{cGammaU}{blue!15}\colorlet{cGammaL}{red!15}\colorlet{dGamma}{black}\colorlet{tGamma}{black}
  \colorlet{cZetaU}{blue!15}\colorlet{cZetaL}{red!15}\colorlet{tZeta}{black}
  \colorlet{cXi}{gray!16}\colorlet{dXi}{gray!48}\colorlet{tXi}{gray!52}
  \colorlet{cNT}{gray!16}\colorlet{dNT}{gray!48}\colorlet{tNT}{gray!52}
  \colorlet{cProd}{gray!16}\colorlet{dProd}{gray!48}\colorlet{tProd}{gray!52}
  \colorlet{cRH}{gray!16}\colorlet{dRH}{gray!48}\colorlet{tRH}{gray!52}
  \colorlet{eArith}{blue!70!black}\colorlet{eElem}{violet!80!black}\colorlet{eDef}{black!80}
  \colorlet{ePNT}{brown!70!black}\colorlet{eAnal}{green!50!black}\colorlet{eExpl}{gray!42}
  \colorlet{eExt}{red!70!black}\colorlet{eZero}{gray!42}\colorlet{eLate}{gray!42}
  \colorlet{lgPrime}{blue!12}\colorlet{lgMu}{teal!15}\colorlet{lgLi}{yellow!25}
  \colorlet{lgXi}{gray!22}\colorlet{lgRH}{gray!22}\colorlet{lgGU}{blue!15}\colorlet{lgGL}{red!15}
  \colorlet{lgEArith}{blue!70!black}\colorlet{lgEDef}{black!80}\colorlet{lgEElem}{violet!80!black}
  \colorlet{lgEZero}{gray!45}\colorlet{lgEAnal}{green!50!black}\colorlet{lgEPNT}{brown!70!black}\colorlet{lgEExt}{red!70!black}
  \colorlet{lgTprime}{black}\colorlet{lgTmu}{black}\colorlet{lgTli}{black}
  \colorlet{lgTxi}{gray!50}\colorlet{lgTrh}{gray!50}\colorlet{lgTgamma}{black}
  \colorlet{lgTarith}{black}\colorlet{lgTdef}{black}\colorlet{lgTelem}{black}
  \colorlet{lgTzero}{gray!50}\colorlet{lgTanal}{black}\colorlet{lgTpnt}{black}\colorlet{lgText}{black}
\else
  \colorlet{cPrime}{blue!12}\colorlet{dPrime}{black}\colorlet{tPrime}{black}
  \colorlet{cMu}{teal!15}\colorlet{dMu}{teal!60!black}\colorlet{tMu}{black}
  \colorlet{cLi}{yellow!25}\colorlet{dLi}{orange!70!black}\colorlet{tLi}{black}
  \colorlet{cGammaU}{blue!15}\colorlet{cGammaL}{red!15}\colorlet{dGamma}{black}\colorlet{tGamma}{black}
  \colorlet{cZetaU}{blue!15}\colorlet{cZetaL}{red!15}\colorlet{tZeta}{black}
  \colorlet{cXi}{cyan!15}\colorlet{dXi}{cyan!60!black}\colorlet{tXi}{black}
  \colorlet{cNT}{cyan!15}\colorlet{dNT}{cyan!60!black}\colorlet{tNT}{black}
  \colorlet{cProd}{cyan!15}\colorlet{dProd}{cyan!60!black}\colorlet{tProd}{black}
  \colorlet{cRH}{orange!15}\colorlet{dRH}{orange!70!black}\colorlet{tRH}{black}
  \colorlet{eArith}{blue!70!black}\colorlet{eElem}{violet!80!black}\colorlet{eDef}{black!80}
  \colorlet{ePNT}{brown!70!black}\colorlet{eAnal}{green!50!black}\colorlet{eExpl}{green!50!black}
  \colorlet{eExt}{red!70!black}\colorlet{eZero}{orange!80!black}\colorlet{eLate}{black!80}
  \colorlet{lgPrime}{blue!12}\colorlet{lgMu}{teal!15}\colorlet{lgLi}{yellow!25}
  \colorlet{lgXi}{cyan!15}\colorlet{lgRH}{orange!15}\colorlet{lgGU}{blue!15}\colorlet{lgGL}{red!15}
  \colorlet{lgEArith}{blue!70!black}\colorlet{lgEDef}{black!80}\colorlet{lgEElem}{violet!80!black}
  \colorlet{lgEZero}{orange!80!black}\colorlet{lgEAnal}{green!50!black}\colorlet{lgEPNT}{brown!70!black}\colorlet{lgEExt}{red!70!black}
  \colorlet{lgTprime}{black}\colorlet{lgTmu}{black}\colorlet{lgTli}{black}
  \colorlet{lgTxi}{black}\colorlet{lgTrh}{black}\colorlet{lgTgamma}{black}
  \colorlet{lgTarith}{black}\colorlet{lgTdef}{black}\colorlet{lgTelem}{black}
  \colorlet{lgTzero}{black}\colorlet{lgTanal}{black}\colorlet{lgTpnt}{black}\colorlet{lgText}{black}
\fi\fi
\makeatother

\begin{tikzpicture}[scale=1.0, transform shape,
  box/.style    = {draw, rounded corners=5pt, minimum width=2.8cm,
                   minimum height=0.9cm, align=center, font=\small, thick},
  zbox/.style   = {draw, rounded corners=5pt, minimum width=2.8cm,
                   minimum height=0.9cm, align=center, font=\small,
                   thick, fill=gray!18},
  splitbox/.style = {draw, rounded corners=5pt, minimum width=7cm,
                     minimum height=2.2cm, align=center, font=\small, thick,
                     rectangle split, rectangle split parts=2,
                     rectangle split part fill={red!8, white}},
  arr/.style    = {-{Stealth[length=6pt,width=5pt]}, thick},
  barr/.style   = {{Stealth[length=6pt,width=5pt]}-{Stealth[length=6pt,width=5pt]}, thick},
  % 线型：虚线仅用于 N(T)→RH、ξ↔RH；其余为实线
  proven/.style = {solid},
  cond/.style   = {dashed},
  lab/.style    = {font=\tiny, align=center, inner sep=3pt},
  rhbox/.style  = {draw=dRH, rounded corners=5pt, minimum width=8cm,
                   minimum height=1.0cm, align=center, font=\small,
                   fill=cRH, text=tRH, thick}
]

% 布局与 A-2 同步：左移 1.0cm、上移 2.0cm
% 左右略放宽，避免 Γ/ζ 自环旁公式标注被裁切
\path[use as bounding box]
  (-10.2, -26.2) rectangle (10.0, 7.4);

\begin{scope}[shift={(-1.0,2.0)}]

\node[font=\bfseries\Large, align=center] at (0, 4.35)
  {黎曼猜想知识谱系关系图};

% 主图单独上移 1cm；标题与图例位置不动
\begin{scope}[shift={(-2.75,-1.2)}]
%% ==================== 节点分层 ====================

% 第一层：素数计数函数 (y=0)
\node[box, fill=cPrime, draw=dPrime]    (J)    at (-4.2,  0.0) {$J(x)$};
\node[box, fill=cPrime, draw=dPrime]    (pi)   at ( 0.0,  0.0) {$\pi(x)$};
\node[box, fill=cPrime, draw=dPrime]    (th)   at ( 5.0,  0.0) {$\theta(x)$};
\node[box, fill=cPrime, draw=dPrime]    (ps)   at ( 9.7,  0.0) {$\psi(x)$};

% 第二层：积性函数 (y=-3.0)
% 节点定义修改
\node[box, fill=cMu, draw=dMu] (mu)  at (0.3, -3.0) {$\mu(n)$};
\node[box, fill=cMu, draw=dMu] (Lam) at (5.0, -3.0) {$\Lambda(n)$};

%Weierstrass乘积
\node[box, fill=cProd, draw=dProd, text=tProd]  (prod)   at ( 9.5,  -17.0) {\tiny\text{Hadamard 乘积}};
% M(x) 节点已移除：初等等价信息直接标注在 RH 框内

% 第三层：Li 与 Gamma (y=-6.0 至 -7.0)
\node[box, fill=cLi, draw=dLi] (Li) at (-2.0, -6.0) {$\Li(x)$};

% Gamma(s) 上下分色：上部标收敛域，下部标延拓
\node[splitbox, rectangle split draw splits=false, minimum width=5cm,
      minimum height=4.7cm,
      rectangle split part fill={cGammaU, cGammaL}] (gamma) at (7.0, -6.0)
    {
     \centering \color{tGamma}$\operatorname{Re}(s)>0$
     \nodepart{two}
     \raisebox{6pt}{\color{tGamma}$\mathbb{C}\setminus\{0,-1,-2,\dots\}$}
    };
% Gamma(s) 标注（固定在节点中线左侧）
\node[font=\small, left=-30pt, text=tGamma] at (gamma.west) {$\Gamma(s)$};

% 【修订3】ζ(s) 节点：下半部分补全函数方程左侧 ζ(s)=
\node[splitbox, rectangle split draw splits=false, minimum height=4cm,
      rectangle split part fill={cZetaU, cZetaL}] (zeta) at (0.3, -10.5)
    {\centering\small
     \color{tZeta}$\displaystyle
       \sum_{n=1}^{\infty}\frac{1}{n^s}
       =\prod_{p}\dfrac{1}{1-p^{-s}}
       \quad(\operatorname{Re}(s)>1)$%
     \vphantom{$\displaystyle\sum_{n=1}^{\infty}$}%
     \nodepart{two}
     % 修订：补全 ζ(s)= 使函数方程完整
     \textcolor{tZeta}{$\displaystyle\zeta(s)=
       2^s\pi^{s-1}
       \sin\!\left(\frac{\pi s}{2}\right)\Gamma(1-s)\,\zeta(1-s)$%
     \vphantom{$\displaystyle\sum_{n=1}^{\infty}\mathbb{C}\setminus\{1\}$}}
    };
% ζ(s) 标注


% ξ(s) 节点：下移至与 RH 同一水平，右侧对称于 M(x)
\node[box, fill=cXi, draw=dXi, text=tXi, minimum height=1.2cm]
  (xi) at (4.85, -17.0)
  {$\xi(s)$\\ \tiny 非平凡零点$\rho$与$\zeta(s)$相同\\
  \tiny $\xi(s)=\xi(1-s)$(函数方程推论)
  };


% 【修订8】N(T) 节点：置于 ζ 与 RH 正中间，构成竖直逻辑链
\node[box, fill=cNT, draw=dNT, text=tNT, minimum height=1.3cm]
  (NT) at (2.3, -14.0)
  {$N(T)$\\ \tiny $\zeta(s)$ 零点计数函数\\[1pt]\tiny（幅角原理）};

% ==================== 黎曼猜想节点 ====================
% RH 下移至 y=-20.5，为 N(T) 留出充分空间
\node[rhbox, minimum width=4.6cm, minimum height=2.2cm] (rh) at (-1.9, -17.0) {
    {\tiny\color{tRH} 黎曼猜想 (Riemann Hypothesis, 1859)}\\[3pt]
    {\tiny\color{tRH}$\zeta(s)$ 的所有非平凡零点 $\rho$ 满足
    $\operatorname{Re}(\rho)=\dfrac{1}{2}$}\\[4pt]
{\tiny\color{tRH} 初等等价（梅滕斯函数）：%
  $\displaystyle M(x)=\sum_{n\leq x}\mu(n)$}\\
{\tiny\color{tRH}
  $\mathrm{RH} \iff M(x) = O(x^{1/2+\varepsilon}),\ \forall \varepsilon>0$}
};


% ==================== 连线 ====================

% ----- 第一层内部 -----
\draw[barr, eArith] (J.east) -- (pi.west)
  node[lab, midway, above=14pt, xshift=-2pt, eArith]
    {$\displaystyle\pi(x)=\sum_{n=1}^{\infty}\frac{\mu(n)}{n}J(x^{1/n})$}
  node[lab, midway, below=14pt, xshift=-2pt, eArith]
    {$\displaystyle J(x)=\sum_{n=1}^{\infty}\frac{1}{n}\pi(x^{1/n})$};

\draw[barr, eElem] (pi.east) -- (th.west)
  node[lab, midway, above=14pt, eElem]
    {$\displaystyle\theta(x)=\pi(x)\log x-\int_2^x\frac{\pi(t)}{t}\,\mathrm{d}t$}
  node[lab, midway, above={14pt}, yshift={-50pt}, eElem]
    {\tiny $\displaystyle\pi(x) = \dfrac{\theta(x)}{\log x}+ \int_{2}^{x} \frac{\theta(t)}{t\,(\log t)^2} \,\mathrm{d}t$};

\draw[barr, eElem] (th.east) -- (ps.west)
  node[lab, midway, above=11pt, xshift={5pt}, eElem]
    {$\displaystyle\psi(x)=\sum_{k=1}^{\infty}\theta(x^{1/k})$}
  node[lab, midway, above=-31pt, xshift={5pt}, yshift={-8pt}, eElem]
    {\tiny $\displaystyle\theta(x)=\sum_{k=1}^{\infty}\mu(k)\psi(x^{1/k})$};


% ----- 第二层内部 -----
\draw[barr, eArith] (mu.east) -- (Lam.west)
  node[lab, midway, yshift=20, xshift=0.3cm, eArith]
    {$\displaystyle\Lambda(n)=-\sum_{d\mid n}\mu(d)\log d$
     \quad{\rm (M\"{o}bius 反演)}}
  node[lab, midway, yshift={-12pt-0.5cm}, xshift=0, eArith]
    {$\displaystyle\log n=\sum_{d\mid n}\Lambda(d)$};

% μ(n) → M(x) 的定义连线已省去（见节点内文字），避免与已有连线交叉

% Lambda -> psi
\draw[arr, eDef] (Lam.east) to[out=0, in=-90, looseness=1.0]
  ($(ps.south west)+(30pt,0)$)
  node[lab, pos=0.5, yshift=-90, xshift=250, eDef]
    {$\displaystyle\psi(x)=\sum_{n\le x}\Lambda(n)$};

% ----- 第三层：Li 和 Gamma -----

% 【修订5】Li(x) → π(x)：改为渐近等价语义（非因果推导）
\draw[arr, ePNT]
  ($(Li.north)!.75!(Li.north east)$)
  to[out=90, in=-135, looseness=1.8]
  ($(pi.south west)+(35pt,0)$)
  node[lab, pos=0.4, yshift=-45, xshift=45, ePNT, align=left]
    {\textbf{素数定理（PNT）}：$\pi(x)\sim\Li(x)$\\
     (若RH成立，$\pi(x)=\Li(x)+O(\sqrt{x}\log x)$)};

% Li -> J（显式公式主项；属 eExpl：第2/9章灰白，第15章全文总图着色）
\draw[arr, eExpl]
  (Li.north) to[out=135, in=-90, looseness=1.5]
  ($(J.south east)+(-22pt,0)$)
  node[lab, pos=0.5, yshift=-115, xshift=-65, eExpl, align=center]
    {$\Li(x)$ 是 $J(x)$ 显式\\[2pt]公式的主项};

% Gamma（Re>0）→ ζ：Euler 积分表示（该积分公式本身要求 Re(s)>1）
\draw[arr, eAnal]
  ([yshift=6pt] gamma.west)
  to[out=180, in=90, looseness=0.8]
  ($(zeta.north)+(30pt,0)$)
  node[lab, pos=0.5, yshift=-235, xshift=90,eAnal, align=center]
    {$\displaystyle\zeta(s)\Gamma(s)
      =\int_0^{\infty}\frac{x^{s-1}}{e^x-1}\,\mathrm{d}x$\\
     $(\operatorname{Re}(s)>1)$};

% μ(n) → ζ(s)
\draw[arr, eAnal]
  (mu.south) to[out=-90, in=90, looseness=1.2]
  ($(zeta.north)+({-20pt+0.3cm},0)$)
  node[lab, pos=0.5, yshift={-120pt-0.5cm-0.4cm}, xshift={45pt-5pt}, eAnal]
    {$\displaystyle\frac{1}{\zeta(s)}=\sum_{n=1}^{\infty}\frac{\mu(n)}{n^s}$\\
     $(\operatorname{Re}(s)>1)$};

% ----- J(x) 与 ψ(x) 的 Stieltjes 积分互逆关系 -----
% 【修订4】加注 dψ(t)/ln t 在 t=1 附近需 t>1 的说明
\draw[barr, eArith, line width=0.8pt, rounded corners=8pt]
  ($(J.north)+(0,0.0cm)$)
  -- ++(0,2.0cm)
  -- ++(13.5cm,0)
  -- ++(0,-2.0cm)
  node[lab, pos=0.35, above=6pt, yshift={10pt+0.5cm}, xshift=-200, eArith, align=center]
    {$\displaystyle\psi(x)=\int_{2^-}^x\log t\,\mathrm{d}J(t)$\quad(Stieltjes 积分，$J\to\psi$)}
  node[lab, pos=0.65, below=6pt, yshift={32pt+0.2cm}, xshift=-205, eArith, align=center]
    {$\displaystyle J(x)=\int_{2^-}^x\frac{\mathrm{d}\psi(t)}{\log t}$
     \quad(Stieltjes 反演，$\psi\to J$；积分自 $t>1$)};

% ----- 第四层：ζ 和 ξ 与第二层的连线 -----

% ζ → ψ：显式公式（Perron 积分 + 留数，单向；x≠p^k 时公式精确成立）
\draw[arr, eExpl, rounded corners=30pt]
  ($(zeta.east)+(0, -10pt)$) -| ($(ps.south)+(30pt,0)$)
  node[lab, pos=0.85, right=10pt, yshift=-170, xshift={-170pt-0.2cm},eExpl, align=left]
{$\psi(x)$ 显式公式（Perron 积分，$\zeta\to\psi$）\\
     $\displaystyle \psi(x) = \frac{1}{2\pi i}\int_{c-i\infty}^{c+i\infty}\!\left(-\frac{\zeta'(s)}{\zeta(s)}\right)\!\frac{x^s}{s}\,\mathrm{d}s$\\
     $\displaystyle = x-\sum_{\rho}\frac{x^{\rho}}{\rho}-\log(2\pi)-\tfrac{1}{2}\log(1-x^{-2})$\\
     \tiny $x>1,\;x\neq p^k$；$\rho$ 按 $|\mathrm{Im}(\rho)|$ 升序配对，条件收敛};

% ψ → ζ：Mellin 变换（单向）
\draw[arr, eExpl, rounded corners=30pt]
  ($(ps.south)+(40pt,0)$) |- ($(zeta.east)+(0, -20pt)$)
  node[lab, pos=0.8, yshift=-27, xshift=30,eExpl, align=center]
    {Mellin 变换（$\psi(x) \longleftrightarrow -\dfrac{\zeta'(s)}{\zeta(s)}$，即负对数导数）\\
     $\displaystyle -\frac{\zeta'(s)}{\zeta(s)} = s\int_{1}^{\infty}\frac{\psi(x)}{x^{s+1}}\,\mathrm{d}x$};
       
       
       
       
       
       
       
% Lambda -> zeta（对数微分 Dirichlet 级数）
\draw[arr, eAnal]
  ($(Lam.south west)+(6pt,0)$) to[out=-90, in=90, looseness=0.4]
  ($(zeta.north west)!0.50!(zeta.north east)$)
  node[lab, pos=0.5, yshift={-120pt-0.5cm-0.4cm+1cm}, xshift={145pt-5pt}, eAnal]
    {$\displaystyle-\frac{\zeta'(s)}{\zeta(s)}
      =\sum_{n=1}^{\infty}\frac{\Lambda(n)}{n^s}$};

% J(x) 显式公式（ζ → J）
\draw[arr, eExpl, rounded corners=50pt]
  ($(zeta.south west)+(0,6pt)$) -- ++(-60pt,0) -- ($(J.south west)+(5pt,0)$)
  node[lab, pos=0.5, yshift=-203, xshift={75pt-0.2cm}, eExpl, align=center]
{\textbf{Perron 公式} $\Rightarrow$ $J(x)$ 显式展开\\
     $\displaystyle
     \begin{aligned}
         &J(x)=\frac{1}{2\pi i}\int_{c-i\infty}^{c+i\infty}
           \frac{\log\zeta(s)}{s}\,x^s\,\mathrm{d}s\\
         &=\Li(x)-\sum_{\rho}\Li(x^{\rho})
           -\log 2+\int_x^{\infty}\frac{\mathrm{d}t}{t(t^2-1)\log t}
     \end{aligned}$\\
     \tiny $x>1,\;x\neq p^k$；$\rho$ 按 $|\mathrm{Im}(\rho)|$ 升序配对，条件收敛};

% ζ → J：Mellin 变换（ln ζ）
\draw[arr, eAnal, rounded corners=5pt]
  (zeta.north west) to[out=110, in=-90, looseness=1.5] (J.south)
  node[lab, pos=0.4, yshift=-205, xshift=-55, eAnal, align=center]
    {$\displaystyle\log\zeta(s)=s\int_1^{\infty}J(x)\,x^{-s-1}\,\mathrm{d}x$\\
     $(\operatorname{Re}(s)>1)$};

% ζ → ξ（定义）：ξ 已下移，连线沿右侧空白区域向右下延伸
\draw[arr, eLate, rounded corners=15pt]
  ($(zeta.south east)!0.15!(zeta.north east)$)
  -- ++(1.0, 0)
  -- ++(0, -3.8)
  -- ($(xi.north)+(0pt,0)$)
  node[lab, pos=0.25, right=4pt,yshift=30, xshift=0, eLate, align=center]
    {$\xi(s)=\dfrac{1}{2}s(s-1)\pi^{-s/2}\Gamma(s/2)\zeta(s)$  
};






% xi -> prod (Weierstrass product connection)
\draw[arr, eLate]
  %(xi.south) -- (prod.south)
  (xi.east) --(prod.west)
  node[midway, right=4pt, yshift=28, xshift=-45, eLate]
    {\tiny$\displaystyle\xi(s)=e^{A+Bs}\prod_{\rho}\left(1-\frac{s}{\rho}\right)e^{s/\rho}$};

% prod -> ps (zero sum in explicit formula)
%\draw[arr, violet!70!black, rounded corners=15pt]
  %(prod.east) -- ++(4.4, 0) |- (ps.south)
  %node[lab, pos=0.6, right=2pt, violet!70!black, align=left]
   % {对数微分与 Perron 积分\\
    % 给出显式公式中\\
    % 零点和 $\sum_{\rho}\frac{x_{\rho}}{\rho}$ 项};
%
% ζ 自身左侧弧线（解析延拓）：与表 tab:spectrum-stage8-guide 一致
\draw[arr, eExt, rounded corners=5pt]
  ($(zeta.north west)!0.15!(zeta.south west)$)
  to[out=180, in=180, looseness=1.6]
  ($(zeta.north west)!0.85!(zeta.south west)$);
% 标签用独立节点锚定在 ζ 西侧弧线旁（避免 path 末 node 漂移）
% 单行；中心在 zeta.west 右 2.5cm（相对初版再累计右移）
\node[lab, eExt, align=center, anchor=center]
  at ([xshift=2.5cm]zeta.west)
  {$\zeta$ 亚纯延拓至 $\mathbb{C}\setminus\{1\}$（Hankel / $\vartheta$ 两法）};

% ζ → N(T)：竖直向下（弧线）
\draw[arr, eZero]
  ($(zeta.south east)!0.10!(zeta.north east)$)
  to[out=0, in=90] ++(0.8, -0.8)
  to[out=-90, in=0] (NT.east)
  node[lab, pos=0.6, right=6pt, yshift=-357, xshift=5, eZero, align=center]
    {$\displaystyle N(T)=\frac{T}{2\pi}\log\frac{T}{2\pi e}+O(\log T)$\\
     \tiny 零点计数公式（von Mangoldt）};

% 【已删去】ζ → RH 蓝色直连线（设计意图改为ζ→N(T)→RH线性链）
%\draw[arr, eArith, rounded corners=10pt]
%  (zeta.south)
%  -- ++(0, -3.35)
%  -- ++(-4.1, 0)
%  -- ($(rh.north)!0.80!(rh.north west)$)
%  node[lab, pos=0.35, left=6pt, yshift=20pt,xshift=120,eArith, align=center]
%    {
%     $\zeta(s)$ 的非平凡零点全部位于直线 $\operatorname{Re}(s)=\dfrac12$ 上};
% RH -- P(s) 画线 + 文字
%\draw[arr, cyan!60!black]
 % ($(rh.east)!0.60!(rh.south east)$) -- (prod.west)    % 起点 -- 终点（必须有！）
  %node[midway, right=6pt, font=\tiny, cyan!60!black, align=center]
  %  {$\zeta(s) = \frac{e^{s(\log(2\pi) - 1)}}{2(s-1)} \prod_{n=1}^{\infty} \left(1 + \frac{s}{2n}\right)e^{-s/2n} P(s)$};

% ξ(s) ↔ RH、N(T) → RH：虚线标注（与 RH 相关的谱系连线）

% N(T) → RH：完成 ζ→N(T)→RH 线性逻辑链（虚线）
\draw[arr, eZero, cond, rounded corners=8pt]
  (NT.south west)
  to[out=-150, in=60, looseness=0.9]
  ($(rh.north)+(20pt, 0)$)
node[lab, pos=0.55, left=4pt, yshift=-430, xshift=-25, eZero, align=center]
    {$\mathrm{RH} \iff N_0(T)=N(T),\ \forall\,T>0$\\
   \tiny（$N_0(T)$：临界线 $\operatorname{Re}(s)=\dfrac{1}{2}$ 上虚部 $\leq T$ 的零点数）};

\draw[barr, eZero, cond]
  (xi.west) -- ($(rh.east)!0.0!(rh.south east)$)
  node[lab, midway, above=20pt,yshift=3pt,xshift=5, eZero, align=center]
    {$\mathrm{RH}\iff\xi\!\left(\dfrac{1}{2}+it\right)$\\
     的所有零点 $t$ 均为实数};

% M(x)→RH 箭头已删去：等价信息直接写在 RH 框内

% Gamma 自身解析延拓弧线：与表 tab:spectrum-stage8-guide 一致
\draw[arr, eExt, rounded corners=1.5pt]
  ($(gamma.north east)+(0,-5pt)$)
  to[out=0, in=0, looseness=1.6]
  ($(gamma.south east)+(0,5pt)$);
% 标签用独立节点锚定在 Γ 东侧弧线旁（仅递推式；延拓域已在框下半）
% 相对初版 (gamma.east+14pt)：净向左 3.7cm，上移 2pt
\node[lab, eExt, align=center, anchor=west]
  at ($(gamma.east)+({14pt-3.7cm},2pt)$)
  {递推延拓 $\Gamma(s)=\Gamma(s+1)/s$};

% Gamma（延拓后）→ ζ：Hankel 围道
\draw[arr, eExt]
  ($(gamma.south west)+(40pt,0)$)
  to[out=180, in=0, looseness=1.0]
  ($(zeta.south east)!0.45!(zeta.north east)$)
  node[lab, pos=0.5, yshift=-210, xshift=215, eExt, align=center]
    {$\displaystyle\zeta(s) = -\frac{\Gamma(1-s)}{2\pi i}\oint_{\mathrm{H}}\frac{(-z)^{s-1}}{e^z-1}\,\mathrm{d}z$\\
     解析延拓到 $\mathbb{C}\setminus\{1\}$，Hankel 围道消去分支切割};

\end{scope}% 主图

% ==================== 图例（取自 A-2 详细文案 + 分阶段灰白色） ====================
\node[draw=black!40, fill=gray!8, rounded corners=8pt, thick, align=center, inner sep=10pt, font=\tiny]
  (legend) at (0, -23.5) {
  \begin{tabular}{@{}ll@{\hspace{1.0cm}}ll@{}}
    \multicolumn{2}{l}{\scriptsize\textbf{节点颜色分类图例}} &
    \multicolumn{2}{l}{\scriptsize\textbf{知识谱系连线语义图例}} \\[6pt]
    \tikz[baseline=-0.5ex]\fill[lgPrime, draw=black!70, thick, rounded corners=1.5pt] (0,-0.12) rectangle (0.55,0.12); &
    \textcolor{lgTprime}{\textbf{素数计数函数}（实变，如 $\pi(x)$）} &
    \textcolor{lgEArith}{\rule[0.5ex]{16pt}{1.4pt}} &
    \textcolor{lgTarith}{\textbf{算术互逆}（M\"{o}bius / Stieltjes 反演）} \\[3pt]
    \tikz[baseline=-0.5ex]\fill[lgMu, draw=black!70, thick, rounded corners=1.5pt] (0,-0.12) rectangle (0.55,0.12); &
    \textcolor{lgTmu}{\textbf{算术积性函数}（离散，如 $\mu(n)$）} &
    \textcolor{lgEDef}{\rule[0.5ex]{16pt}{1.4pt}} &
    \textcolor{lgTdef}{\textbf{对象定义}（构造定义与主项提取）} \\[3pt]
    \tikz[baseline=-0.5ex]\fill[lgLi, draw=black!70, thick, rounded corners=1.5pt] (0,-0.12) rectangle (0.55,0.12); &
    \textcolor{lgTli}{\textbf{渐近逼近主项}（如 $\Li(x)$）} &
    \textcolor{lgEElem}{\rule[0.5ex]{16pt}{1.4pt}} &
    \textcolor{lgTelem}{\textbf{初等计数变换}（数论函数基本恒等式）} \\[3pt]
    \tikz[baseline=-0.5ex]\fill[lgXi, draw=black!55, thick, rounded corners=1.5pt] (0,-0.12) rectangle (0.55,0.12); &
    \textcolor{lgTxi}{\textbf{辅助复变函数}（如 $\xi(s)$）} &
    \textcolor{lgEZero}{\rule[0.5ex]{16pt}{1.4pt}} &
    \textcolor{lgTzero}{\textbf{零点结构与猜想}（分布特征及等价命题）} \\[3pt]
    \tikz[baseline=-0.5ex]\fill[lgRH, draw=black!55, thick, rounded corners=1.5pt] (0,-0.12) rectangle (0.55,0.12); &
    \textcolor{lgTrh}{\textbf{核心猜想与目标}（如 RH）} &
    \textcolor{lgEAnal}{\rule[0.5ex]{16pt}{1.4pt}} &
    \textcolor{lgTanal}{\textbf{分析桥梁}（积分变换与 Dirichlet 级数）} \\[3pt]
    \tikz[baseline=-0.5ex]{\clip[rounded corners=1.5pt] (0,-0.12) rectangle (0.55,0.12); \fill[lgGL] (0,-0.12) rectangle (0.55,0); \fill[lgGU] (0,0) rectangle (0.55,0.12); \draw[black!55, thick, rounded corners=1.5pt] (0,-0.12) rectangle (0.55,0.12); \draw[black!55, thick] (0,0) -- (0.55,0);} &
    \textcolor{lgTgamma}{\textbf{解析延拓复变}（初始域/延拓域）} &
    \textcolor{lgEPNT}{\rule[0.5ex]{16pt}{1.4pt}} &
    \textcolor{lgTpnt}{\textbf{渐近定理}（PNT 等确立的渐近推论）} \\[3pt]
    & & \textcolor{lgEExt}{\rule[0.5ex]{16pt}{1.4pt}} &
    \textcolor{lgText}{\textbf{解析延拓}（复围道积分与全纯延拓）} \\[3pt]
    & & \tikz[baseline=-0.5ex]{\draw[lgEZero, line width=1.2pt] (0.00,0) -- (0.12,0); \draw[lgEZero, line width=1.2pt] (0.17,0) -- (0.29,0); \draw[lgEZero, line width=1.2pt] (0.34,0) -- (0.46,0);} &
    \textbf{未证}
  \end{tabular}
};

\end{scope}% 位移：左 1.0cm、上 2.0cm

\end{tikzpicture}
%
}
\end{figure}
\clearpage
\noindent
\captionof{figure}{黎曼猜想知识谱系关系图（第一次着色：初等算术层）。
与图~\ref{fig:function-relations} 同一底图；彩色为本章已建立的主要对象与关系，灰白为后文内容。
彩色箭头与本章正文的逐项对照见表~\ref{tab:spectrum-stage2-guide}。}
\label{fig:spectrum-stage2}

\paragraph{图~\ref{fig:spectrum-stage2} 彩色元素与本章正文对照}
下列各项将图上\textbf{已着色}的节点与连线引向本章相应定义与证明；
灰白部分待第二次着色（第~\ref{ch:riemann_zeta_continuation}~章）
与第三次全文着色（第~\ref{ch:spectrum_map}~章）后再逐条对接正文。

\begin{table}[htbp]
\centering
\caption{图~\ref{fig:spectrum-stage2} 彩色图元与第~\ref{ch:elementary_number_theoretic_functions}~章正文对照}
\label{tab:spectrum-stage2-guide}
\small
\renewcommand{\arraystretch}{1.35}
\begin{tabular}{@{}p{4.8cm}p{8.4cm}@{}}
\toprule
\textbf{图上元素} & \textbf{本章正文（定义 / 公式 / 证明）} \\
\midrule
节点 $\pi(x)$ &
  第~\ref{sec:prime_counting}~节：定义、素数定理陈述与数值 \\
节点 $\Li(x)$ &
  第~\ref{sec:offset_logarithmic_integral}~节：$\Li$ 从简；完整 $\li$/$\Ei$ 见第~\ref{ch:logint_Ei}~章 \\
连线 $\Li(x)\to\pi(x)$（PNT） &
  第~\ref{subsec:Li_pi_asymp}~节：$\pi\sim\Li$ 及大尺度数值 \\
节点 $J(x)$ &
  第~\ref{sec:J}~节：定义 \eqref{eq:J_def}（$J=\sum_{p^k\le x}1/k$） \\
连线 $J(x)\leftrightarrow\pi(x)$
  （Möbius 型互换） &
  第~\ref{subsec:J_from_pi}~、\ref{subsec:J_pi_mobius}~节：
  \eqref{eq:J_from_pi}、\eqref{eq:pi_from_J}
  （求和至 $N(x)=\lfloor\log x/\log 2\rfloor$） \\
节点 $\mu(n)$ 及 RH 框内 $M(x)$ 信息 &
  第~\ref{sec:mobius_function}~节：定义与反演；
  第~\ref{subsec:mertens_function}~节：$M(x)$ 定义、
  RH 含义一句与 $M$ 的初等等价陈述；
  生成函数见第~\ref{ch:riemann_zeta_intro}~章；
  等价性证明见第~\ref{ch:riemann_explicit_formula}~章 \\
节点 $\Lambda(n)$ &
  第~\ref{sec:Lambda}~节：定义与 \eqref{eq:log_sum_Lambda} \\
连线 $\mu\leftrightarrow\Lambda$（卷积 / 反演） &
  第~\ref{subsec:Lambda_mu_conv}~节：
  \eqref{eq:Lambda_mu_conv}、\eqref{eq:Lambda_mu_log}；
  $-\zeta'/\zeta=\sum\Lambda n^{-s}$ 见第~\ref{ch:riemann_zeta_intro}~章 \\
连线 $\Lambda\to\psi$（定义求和） &
  第~\ref{sec:second_chebyshev}~节：$\psi(x)=\sum_{n\le x}\Lambda(n)$ \\
节点 $\theta(x)$、$\psi(x)$ &
  第~\ref{sec:first_chebyshev}~、\ref{sec:second_chebyshev}~节；
  $\psi$ 定义 \eqref{eq:psi_def} \\
连线 $\pi\leftrightarrow\theta$、$\theta\leftrightarrow\psi$ &
  第~\ref{subsec:theta_pi_abel}~、\ref{subsec:theta_psi_PNT}~、
  \ref{subsec:psi_theta_rel}~节：
  分部求和 \eqref{eq:theta_pi_abel}；
  \eqref{eq:psi_sum_theta}、\eqref{eq:psi_theta_asymp}、\eqref{eq:theta_from_psi}；
  $\pi\sim x/\log x$、$\theta\sim x$、$\psi\sim x$ 等价（证明提要） \\
连线 $J\leftrightarrow\psi$（积分 / 阶梯 Stieltjes） &
  第~\ref{subsec:J_psi_stieltjes}~节；
  定义见附录~\ref{app:abel_stieltjes}；
  \eqref{eq:psi_from_J_stieltjes}、\eqref{eq:J_from_psi_stieltjes}、\eqref{eq:psi_J_parts} \\
灰白节点 $\zeta,\Gamma,\xi,N(T)$、RH 框主体
  及灰白连线/公式 &
  本章未建立其解析内容；见第~\ref{ch:riemann_zeta_continuation}~章
  图~\ref{fig:spectrum-stage8} 与第~\ref{ch:spectrum_map}~章
  图~\ref{fig:function-relations} 的后续着色与对照 \\
\bottomrule
\end{tabular}
\end{table}

后文在同一底图上继续着色：第二次（第~\ref{ch:riemann_zeta_continuation}~章末）叠加 $\zeta$ 与函数方程相关部分，
第三次（第~\ref{ch:spectrum_map}~章）给出全文总图，并同样用「图元 $\to$ 正文小节」索引证明。

\section*{本章小结与后续展望}\label{sec:summary_elementary_functions}
\addcontentsline{toc}{section}{本章小结与后续展望}
 
本章给出后文所需的初等对象：$\pi$、$\Li$、$J$、$\mu$、$M$、$\Lambda$、$\theta$、$\psi$。它们定义在实数或正整数上，彼此由反演、分部求和与阶梯积分联系（Abel / Stieltjes 的系统表述见附录~\ref{app:abel_stieltjes}）。与 $\zeta$ 的 Dirichlet 级数关系（$\mu$、$\Lambda$ 为系数）在第~\ref{ch:riemann_zeta_intro}~章于 $\operatorname{Re}s>1$ 上建立；解析延拓与显式公式留待后文。
 
\begin{center}
\fbox{\parbox{0.92\textwidth}{
\centering
\vspace{0.2em}
\textbf{第\ref{ch:elementary_number_theoretic_functions}章：重要定理与核心公式汇总}\\[6pt]
\renewcommand{\arraystretch}{1.6}
\begin{tabular}{lp{9.5cm}}
\hline
\textbf{名称} & \textbf{数学内容 / 公式表达} \\
\hline
Möbius 反演公式 & $g(n) = \displaystyle\sum_{d|n} f(d) \iff f(n) = \displaystyle\sum_{d|n} \mu(d)g(n/d)$ \\
Chebyshev 与 PNT 等价 &
  $\theta=\pi\log x-\displaystyle\int_2^x\pi(t)/t\,\mathrm{d}t$~\eqref{eq:theta_pi_abel}；
  $\psi=\displaystyle\sum_k\theta(x^{1/k})=\theta+O(\sqrt{x}\log x)$~\eqref{eq:psi_sum_theta}--\eqref{eq:psi_theta_asymp}；
  $\theta=\displaystyle\sum\mu(k)\psi(x^{1/k})$~\eqref{eq:theta_from_psi}；
  $\pi\sim x/\log x\iff\theta\sim x\iff\psi\sim x$（证明提要） \\
$\psi(x)$ 定义 &
  $\psi=\displaystyle\sum_{p^k\le x}\log p=\displaystyle\sum_{n\le x}\Lambda(n)$~\eqref{eq:psi_def} \\
$J(x)$ 的定义与 $\pi$ 形 & 定义 $J(x)=\displaystyle\sum_{p^k\le x}1/k$~\eqref{eq:J_def}；
  等价 $J(x)=\displaystyle\sum_{n=1}^{N(x)}\frac{1}{n}\pi(x^{1/n})$~\eqref{eq:J_from_pi}；
  反演 $\pi(x)=\displaystyle\sum_{n=1}^{N(x)}\frac{\mu(n)}{n}J(x^{1/n})$~\eqref{eq:pi_from_J}
  （$N(x)=\lfloor\log x/\log 2\rfloor$） \\
$J(x)$ 与 $\psi(x)$ 积分关系 &
  $\psi=\displaystyle\int_{2^-}^x\log t\,\mathrm{d}J$~\eqref{eq:psi_from_J_stieltjes}，
  $J=\displaystyle\int_{2^-}^x\mathrm{d}\psi/\log t$~\eqref{eq:J_from_psi_stieltjes}，
  分部 \eqref{eq:psi_J_parts}
  （见第~\ref{subsec:J_psi_stieltjes}~节） \\
\hline
\end{tabular}
\vspace{0.2em}
}}
\end{center}
 
第~\ref{ch:riemann_zeta_intro}~章将建立 $\mu$、$\Lambda$ 与 $\zeta$、$1/\zeta$、$-\zeta'/\zeta$ 在 $\operatorname{Re}s>1$ 上的 Dirichlet 级数联系；此后 $J$、$\psi$ 经 Perron 积分与显式公式与非平凡零点相连。下一章准备复分析工具。